\documentclass[UTF8,openany,zihao=-4]{ctexbook}

% Basic page layout
\usepackage[a4paper,margin=2.6cm,headheight=15pt]{geometry}
\usepackage{setspace}
\onehalfspacing
\usepackage{amsmath}

% Figures and tables
\usepackage{graphicx}
\graphicspath{{figures/}}
\usepackage{booktabs}
\usepackage{tabularx}
\usepackage{array}
\usepackage{caption}
\usepackage{subcaption}

% Lists, colors, and links
\usepackage{xcolor}
\usepackage{enumitem}
\setlist{nosep,leftmargin=2em}
\usepackage{csquotes}
\usepackage[colorlinks=true,linkcolor=blue!55!black,citecolor=blue!55!black,urlcolor=blue!55!black]{hyperref}

% Headings and running headers
\usepackage{titlesec}
\titleformat{\chapter}[hang]{\Huge\bfseries}{第\,\thechapter\,章}{1em}{}
\titlespacing*{\chapter}{0pt}{-0.5em}{1.5em}

\usepackage{fancyhdr}
\pagestyle{fancy}
\fancyhf{}
\fancyhead[LE,RO]{\thepage}
\fancyhead[LO]{\nouppercase{\rightmark}}
\fancyhead[RE]{\nouppercase{\leftmark}}

% Bibliography
\usepackage[backend=biber,style=numeric,sorting=nyt]{biblatex}
\addbibresource{references.bib}

% Small teaching helpers for later drafting
% 手牌/牌面：Latin Modern 与微软雅黑均不含 ♠♣♦♥，需单独映射符号字体
\IfFontExistsTF{Segoe UI Symbol}{%
  \newfontfamily\SuitFont{Segoe UI Symbol}[Scale=0.92]
}{%
  \IfFontExistsTF{DejaVu Sans}{%
    \newfontfamily\SuitFont{DejaVu Sans}[Scale=0.92]
  }{%
    \IfFontExistsTF{Arial Unicode MS}{%
      \newfontfamily\SuitFont{Arial Unicode MS}[Scale=0.92]
    }{%
      \newfontfamily\SuitFont{Symbola}[Scale=0.92]
    }%
  }%
}
\usepackage{newunicodechar}
\newunicodechar{♠}{{\SuitFont ♠}}
\newunicodechar{♥}{{\SuitFont ♥}}
\newunicodechar{♦}{{\SuitFont ♦}}
\newunicodechar{♣}{{\SuitFont ♣}}
\newcommand{\hand}[1]{{\normalfont #1}}
\newcommand{\term}[2]{#1（\textit{#2}）}
\newcommand{\todoitem}[1]{\textcolor{red!70!black}{[待补充：#1]}}


\title{告别凭感觉打牌：德州扑克新手系统入门课}
\author{作者待定}
\date{\today}

\begin{document}

\frontmatter

\begin{titlepage}
  \centering
  \vspace*{3cm}
  {\zihao{1}\bfseries 告别凭感觉打牌\par}
  \vspace{0.8cm}
  {\zihao{2}\bfseries 德州扑克新手系统入门课\par}
  \vspace{1.5cm}
  {\zihao{4} 从直觉决策到框架化思考\par}
  \vfill
  {\zihao{4} 作者待定\par}
  \vspace{0.5cm}
  {\zihao{4}\today\par}
\end{titlepage}

\chapter*{前言}
\addcontentsline{toc}{chapter}{前言}

很多人第一次接触德州扑克，都会被它表面的简单所吸引：每个人两张底牌，五张公共牌，最后比谁的牌型更大。规则似乎不难，几分钟就能学会；但真正坐上牌桌后，才会发现每一手牌都充满了不确定性。

同样是顶对，有时可以大胆价值下注，有时却只是勉强抓诈唬；同样是同花听牌，有时适合主动半诈唬，有时面对大注应该果断弃牌；同样是河牌大注，有时代表极强价值牌，有时却可能是对手错过听牌后的最后一击。很多新手的困惑，并不是不知道牌型大小，而是不知道在具体局面中应该如何思考。

他们常常会问：

\begin{itemize}
\item 这手牌翻前到底能不能玩？
\item 翻牌中了顶对，是不是一定要下注？
\item 我有听牌，能不能继续跟？
\item 转牌变危险了，还要不要开第二枪？
\item 河牌面对大注，我到底该不该跟？
\item 我下注到底是为了价值，还是为了诈唬？
\item 为什么有时候弃牌才是正确选择？
\end{itemize}

这些问题看似分散，本质上都指向同一个核心：你是否拥有一套稳定的决策框架。

德州扑克最容易误导新手的地方在于，短期结果会掩盖决策质量。你可能用错误的跟注抓到一次诈唬，也可能用正确的弃牌错过一次对手 bluff；你可能用糟糕的翻前范围击中两对赢下大底池，也可能用标准打法被河牌反超。短期输赢并不总能告诉你自己打得好不好。如果只凭结果复盘，就很容易把运气当成能力，把偶然当成经验，把错误打法当成“有效技巧”。

本书的目的，不是教你背诵复杂公式，也不是让你一夜之间掌握所有高级理论，而是帮助你完成一个最重要的转变：

\begin{center}
\textbf{从凭感觉打牌，转向有框架、有依据、有复盘标准地打牌。}
\end{center}

所谓“框架”，不是死板答案。扑克中几乎没有一个动作在所有局面下都永远正确。真正有用的框架，是让你在复杂局面中知道该看什么、该问什么、该如何逐步缩小错误空间。

因此，本书不会从一堆抽象术语开始，而是按照一名新手最容易理解的顺序展开。

第一章先建立牌手认知。德州扑克不是单纯的纸牌游戏，而是一项包含概率、信息、心理、资金管理和长期执行能力的综合竞技。想成为更好的牌手，首先要理解体能、心态、技术、复盘和资金管理为什么同样重要。

第二章到第五章，按照一手牌的自然流程展开：翻前、翻牌、转牌、河牌。翻前是整手牌的地基，决定你用什么范围进入底池；翻牌是公共牌第一次出现，要求你理解牌面结构和双方范围；转牌是变化最明显的一条街，很多牌力会因为一张转牌而重新定位；河牌则是最终决策街，没有未来补牌机会，所有行动都要回到价值、诈唬、摊牌或弃牌。

第六章开始从单手牌流程上升到范围视角。扑克不是简单地把牌分成强牌和弱牌。坚果牌、强价值牌、中等摊牌价值、弱摊牌价值、强听牌、弱听牌和空气牌，在整体范围中承担着不同功能。理解这些功能，是从“我有什么牌”升级到“这手牌应该完成什么任务”的关键一步。

第七章讨论下注的目的。很多新手的问题不是不敢下注，而是不知道自己为什么下注。下注的核心目的只有两个：价值和诈唬。价值下注要让更差牌跟注，诈唬下注要让更好牌弃牌。如果两个问题都回答不了，就不应该轻易把筹码推进底池。

第八章讨论激进的游戏。德州扑克不是被动等待好牌的游戏。赢得底池有两种方式：摊牌时拥有最好牌，或者通过进攻让对手弃牌。有效激进不是乱 bluff，而是建立在范围优势、手牌功能和对手可弃牌范围之上的有条件施压。

第九章补充基础计算：Outs、胜率与赔率。它是第十章“弃牌与跟注”的数学基础。你需要明白，继续一手牌不能只看“有没有希望”，而要看这个希望值不值当前价格。数 outs、估胜率、算底池赔率，是新手必须掌握的基础工具。

第十章讨论弃牌与跟注。它是最重要、也最反直觉的一章。很多娱乐玩家最大的漏洞不是诈唬太少，而是跟注太多。跟注不是情绪行为，而是数学和范围判断；弃牌不是软弱，而是纪律。能在错误价格下放弃，能克服“不想被诈唬”的心理，是长期进步的关键。

全书的主线可以概括为：

\begin{center}
\textbf{先建立认知，再掌握流程；\\
先理解牌力功能，再明确下注目的；\\
先学会合理进攻，再学会理性防守。}
\end{center}

在阅读本书时，我建议你不要把每一章当成单独知识点，而要把它们连接成一套完整决策系统。比如，翻前范围会影响翻牌范围优势；翻牌下注目的会影响转牌继续计划；转牌出张会改变河牌价值下注或抓诈唬条件；牌力分类会决定底池大小；底池赔率会决定跟注是否划算。

这也是德州扑克最有魅力的地方：它不是简单地记住“什么牌该打，什么牌不该打”，而是在不断变化的信息中做出长期期望值更高的选择。

需要特别说明的是，本书并不承诺任何人通过学习就一定能盈利。扑克存在波动，也涉及资金风险。任何理论和框架都不能消除短期运气，也不能替代稳定心态、合理资金管理和长期复盘。学习本书的目的，应当是提升理解能力和决策质量，而不是追求短期结果。

如果你是初学者，希望你读完本书后，至少能做到三件事：

第一，少用感觉做决定。面对下注、跟注、加注和弃牌时，先问清楚自己的理由。

第二，少犯昂贵错误。尤其是翻前过宽入池、中等牌力打大底池、听牌无脑追、河牌不舍得弃这些新手高频错误。

第三，建立复盘标准。赢牌不一定打得好，输钱也不一定打错。你要学会从位置、范围、牌面、下注目的、赔率和对手类型出发，判断自己的决策是否合理。

本书的名字可以理解为一句提醒：告别凭感觉打牌。

真正的进步，不是某一天突然学会一个神奇技巧，而是你在一手又一手牌中，逐渐用清晰的问题替代模糊的冲动，用稳定的框架替代短期情绪，用长期决策质量替代一时输赢。

愿这本书能帮助你在牌桌上少一点迷茫，多一点纪律；少一点冲动，多一点依据；少一点“我感觉”，多一点“我知道为什么”。

\tableofcontents

\mainmatter

\chapter{如何成为优秀的牌手}

\section{本章导入}

本章将帮助读者先建立对德州扑克的正确认识：它不是依赖短期结果和情绪冲动的游戏，而是一项围绕概率、决策质量、心理控制和风险管理展开的竞技活动。对新手而言，成为优秀牌手的第一步不是记住更多技巧，而是建立长期主义的学习方式。

优秀牌手通常具备三类基础能力：稳定的体能与专注力、成熟的心态管理，以及可复盘的技术框架。本章后续将围绕这三项基础展开，并说明 EV 思维、资金管理和复盘习惯为什么比单手牌的输赢更重要。

\section{核心概念}

本章重点关注长期期望值（EV）、情绪控制、风险管理、资金管理和复盘流程。读者需要理解，短期结果可能受运气影响，但长期收益更接近决策质量的累积结果。

\section{新手常见误区}

新手常把赢钱等同于打得好，把输钱等同于打得差；也容易在连续失利后进入情绪化决策。本节后续将整理这些误区，并引导读者用过程指标替代结果崇拜。

\section{实战思考框架}

面对一手牌，优秀牌手会先问：我的范围是什么，对手的范围是什么，当前行动的目的是什么，风险是否在可承受范围内。这个框架会贯穿全书后续章节。

\section{牌例拆解}

\subsection{牌例一：待补充}

本牌例将用于展示短期结果与长期正确决策之间的区别。

\subsection{牌例二：待补充}

本牌例将用于展示情绪化追损如何破坏原本合理的策略。

\section{本章总结}

本章强调优秀牌手的基础不是神秘直觉，而是长期主义、稳定心态、风险控制和持续复盘。后续章节将在这一基础上进入具体街次决策。

\section{课后训练}

请记录最近三手印象深刻的牌，并尝试分别写下：当时的决策理由、是否受情绪影响、复盘后是否仍然认可该决策。

\chapter{翻前 / Pre-flop}

\section{本章导入}

本章将围绕翻前决策展开。翻前是整手牌的起点，位置、起手牌范围和主动权会直接影响后续每条街的决策质量。初学者常见的问题不是翻后完全不会打，而是翻前已经用过宽、过弱、过被动的范围进入了困难局面。

翻前学习的目标，是建立清晰的入池纪律：什么位置可以 open raise，哪些牌适合 call，什么时候考虑 3-bet，什么时候应当直接 fold。

\section{核心概念}

本章将介绍位置、起手牌范围（range）、open raise、call、3-bet、fold 和主动权。读者需要理解，翻前不是孤立地评价两张手牌，而是结合位置、对手类型和后续可操作性进行判断。

\section{新手常见误区}

娱乐玩家常见错误包括入池范围过宽、喜欢 limp、面对加注过度跟注、低估位置劣势，以及拿边缘牌制造复杂局面。本节后续将逐项说明这些错误如何放大翻后压力。

\section{实战思考框架}

翻前可以按照“位置--手牌--前序行动--对手倾向--后续计划”的顺序思考。每一次入池都应当有理由，而不是因为“这手牌看起来还可以”。

\section{牌例拆解}

\subsection{牌例一：待补充}

本牌例将用于分析早期位置用边缘牌入池带来的连锁问题。

\subsection{牌例二：待补充}

本牌例将用于比较 call 与 3-bet 在不同对手类型下的策略差异。

\section{本章总结}

翻前质量决定了整手牌的起点。新手应先建立紧凑、清晰、有位置意识的范围，再逐步学习更复杂的调整。

\section{课后训练}

请为 UTG、CO、BTN 三个位置分别列出你愿意率先加注的起手牌范围，并标记其中最容易误玩的边缘牌。

\chapter{翻牌 / Flop}

\section{本章导入}

本章将进入翻牌圈。翻牌不是简单地“中了就打，没中就弃”，而是需要结合牌面结构、双方范围、位置关系和下注目的进行判断。很多新手在翻牌圈犯错，是因为只看自己的手牌强弱，而忽略了整体局面。

翻牌学习的重点，是理解范围优势、坚果优势、持续下注（C-bet）以及不同牌面上的价值下注和诈唬下注逻辑。

\section{核心概念}

本章将介绍牌面结构、干燥牌面、湿润牌面、范围优势（range advantage）、坚果优势（nut advantage）、价值下注、诈唬下注和 C-bet。读者需要学会从“我有什么牌”转向“我的范围在这个牌面上表现如何”。

\section{新手常见误区}

新手常见误区包括任何未成牌都自动放弃、任何顶对都过度保护、所有牌面都机械 C-bet，以及用没有后续计划的诈唬下注制造困难局面。

\section{实战思考框架}

翻牌圈可以先评估牌面是否击中自己的翻前范围，再判断手牌属于价值、摊牌价值、听牌还是空气牌。行动前应明确下注、过牌或跟注的目的。

\section{牌例拆解}

\subsection{牌例一：待补充}

本牌例将用于分析干燥高牌面上的持续下注。

\subsection{牌例二：待补充}

本牌例将用于分析湿润连接牌面上的控池与保护。

\section{本章总结}

翻牌圈是范围思维的第一道关口。读者应避免只凭命中程度行动，而要把牌面、范围和下注目的结合起来。

\section{课后训练}

请选择三个不同翻牌面，分别判断它们更有利于翻前进攻方还是防守方，并写出你的理由。

\chapter{转牌 / Turn}

\section{本章导入}

本章讨论转牌圈。转牌的出现会显著改变局势：一些牌会增强进攻方范围，一些牌会补全防守方听牌，也有一些牌会让原本清晰的价值牌变得尴尬。转牌不是翻牌策略的机械延续，而是一次重新评估。

学习转牌的目标，是理解哪些转牌适合继续进攻，哪些转牌适合控池，哪些局面应当承认计划失败并放弃。

\section{核心概念}

本章将介绍好转牌、坏转牌、范围变化、第二枪、控池、延迟价值下注和放弃策略。读者需要看到，转牌决策往往比翻牌更考验计划性。

\section{新手常见误区}

新手常在转牌犯两类错误：一类是翻牌下注后无论转牌如何都继续下注，另一类是一遇到压力就放弃所有中等牌力。本节后续将讨论如何在二者之间找到依据。

\section{实战思考框架}

转牌应按“转牌是否改变坚果结构--谁的范围更受益--我的手牌功能是否变化--下注后河牌计划是什么”的顺序思考。没有河牌计划的转牌下注，往往会制造更昂贵的错误。

\section{牌例拆解}

\subsection{牌例一：待补充}

本牌例将用于分析进攻方在有利转牌上的第二枪。

\subsection{牌例二：待补充}

本牌例将用于分析防守方在危险转牌上的跟注与弃牌边界。

\section{本章总结}

转牌是局势变化最明显的一街。读者需要学会根据出张重新评估范围和牌力，而不是依赖翻牌圈的第一印象。

\section{课后训练}

请找出三个你曾经在转牌犹豫的牌局，分别判断转牌属于好牌、坏牌还是中性牌，并说明原因。

\chapter{河牌 / River}

\section{本章导入}

本章讨论河牌圈。河牌是最终决策街，所有听牌要么完成，要么落空，双方都不再拥有未来街次的补偿空间。因此，河牌决策必须更加明确：是价值下注、诈唬下注、抓诈唬，还是放弃摊牌。

新手在河牌最容易被结果情绪影响：害怕被反超而不敢价值下注，或因为不甘心而用错误价格跟注。本章将帮助读者建立更稳定的河牌判断标准。

\section{核心概念}

本章将介绍摊牌价值、价值下注、薄价值下注、诈唬下注、抓诈唬、阻断牌和面对大注的决策。读者需要理解，河牌行动的依据来自范围、赔率和对手倾向，而不是单纯的害怕或好奇。

\section{新手常见误区}

常见误区包括错过价值下注、用中等牌力下注后只能被更好牌跟注、面对大注过度好奇跟注，以及在没有弃牌率的对手面前强行诈唬。

\section{实战思考框架}

河牌可以先判断自己的牌是否能被更差牌跟注，再判断对手是否会弃掉更好牌。面对下注时，则应结合底池赔率、对手价值范围和潜在诈唬组合做决定。

\section{牌例拆解}

\subsection{牌例一：待补充}

本牌例将用于分析河牌薄价值下注的边界。

\subsection{牌例二：待补充}

本牌例将用于分析面对河牌大注时的抓诈唬条件。

\section{本章总结}

河牌要求读者用最清晰的标准做最终选择。好的河牌决策不是勇敢或保守，而是价格、范围和对手倾向共同支持的结果。

\section{课后训练}

请复盘三手河牌面对下注的牌，计算当时跟注所需胜率，并判断自己的跟注是否有数学依据。

\chapter{不同牌力的处理框架}

\section{本章导入}

本章从单手牌思维上升到牌力分类框架。扑克是范围的游戏，同一手牌在不同牌面、位置和对手行动下可能承担不同功能。初学者需要先学会给自己的手牌分类，再决定它适合下注、过牌、跟注、加注还是弃牌。

本章将为后续下注目的、激进策略和弃牌跟注打下基础。只有理解不同牌力在整体范围中的角色，才能避免把所有牌都当成“强或弱”这种过于粗糙的二分法。

\section{核心概念}

本章将介绍坚果牌、强价值牌、中等摊牌价值、弱摊牌价值、强听牌、弱听牌和空气牌。读者需要理解，牌力分类不是固定标签，而是随公共牌和行动线变化的动态判断。

\section{新手常见误区}

新手常把顶对自动当成强牌，把听牌自动当成必须追的牌，也常低估中等摊牌价值的过牌意义。本节后续将说明不同牌力的主要任务。

\section{实战思考框架}

每一街行动前，可以先问：我的牌属于价值、摊牌价值、听牌还是空气牌？它需要保护吗？它能从更差牌获得价值吗？它能让更好牌弃牌吗？

\section{牌例拆解}

\subsection{牌例一：待补充}

本牌例将用于展示同一手顶对在不同牌面结构下的牌力变化。

\subsection{牌例二：待补充}

本牌例将用于展示强听牌和弱听牌在进攻策略上的差异。

\section{本章总结}

牌力分类是框架化决策的重要工具。读者应学会先识别手牌功能，再选择对应行动，而不是直接根据情绪决定进退。

\section{课后训练}

请选取五个翻牌局面，把自己的手牌分别归入本章的牌力类别，并写出最自然的行动方案。

\chapter{下注的目的}

\section{本章导入}

下注，是德州扑克中最重要的主动行动之一。你在牌桌上没有语言可以解释自己的想法，也不能直接要求对手弃牌或支付价值。你真正能使用的工具，只有筹码。你把筹码推进底池，对手必须根据你的下注尺度、行动线、牌面结构和自身牌力作出回应。

从这个角度看，筹码就是你唯一的武器，下注就是你最主要的进攻方式。

很多新手并不是不敢下注，而是不知道自己为什么下注。他们会因为“我有牌”而下注，会因为“我怕对手反超”而下注，会因为“我想知道他有没有”而下注，会因为“我不想让他免费看牌”而下注，也会因为“我觉得这里应该打一下”而下注。表面上看，这些理由都像是下注原因，但它们并不够清晰。

真正有效的下注，必须回到两个核心目的：价值和诈唬。

如果你是价值下注，你希望对手用更差的牌跟注。
如果你是诈唬下注，你希望对手弃掉比你更好的牌。

所有其他概念，例如保护、试探、控池、施压、阻止对手免费看牌，都不能脱离这两个核心目的独立存在。保护性下注本质上仍然要么是在向更差牌和听牌收费，要么是在迫使对手放弃有权益的牌；半诈唬本质上是诈唬的一种，但它同时保留了后续改善成强牌的机会；薄价值下注本质上仍然是价值下注，只是它面对的是更窄、更边缘的更差跟注范围。

本章是整门课程的核心之一。前面我们已经学习了翻前、翻牌、转牌、河牌的处理流程，也学习了不同牌力在范围中的功能。从本章开始，我们进一步追问一个更直接的问题：

\begin{center}
\textbf{我下注到底想让对手做什么？}
\end{center}

本章的核心输出是“下注前必问”：

\begin{center}
\textbf{如果我是价值下注，哪些更差牌会跟？\\
如果我是诈唬，哪些更好牌会弃？\\
如果两个问题都回答不了，就不要轻易下注。}
\end{center}

只要你能在每次下注前问这三个问题，就会立刻减少大量无目的下注、情绪下注和试探性下注。

\section{核心概念}

\subsection{筹码是你唯一的武器}

德州扑克是一项信息不完全的竞技。你不知道对手的底牌，对手也不知道你的底牌。你们之间的信息交流，主要通过行动完成：过牌、下注、跟注、加注、弃牌。所有行动中，下注最具主动性，因为下注会迫使对手立刻作出选择。

当你下注时，对手通常只有几种回应方式：弃牌、跟注、加注。你通过下注改变了底池大小，也改变了对手继续游戏所需要付出的成本。

因此，下注不是简单地“把钱放进底池”。它至少有三层意义：

\begin{enumerate}
\item 它表达了你的范围强度；
\item 它改变了对手继续所需的价格；
\item 它迫使对手用自己的范围作出筛选。
\end{enumerate}

例如，你在翻牌下注 1/3 底池，给对手的压力相对较小，对手可以用较宽范围继续。你在转牌下注 2/3 底池，对手继续成本明显提高，需要放弃更多弱牌。你在河牌使用超池下注，对手往往只能用很强的牌或非常合适的抓诈唬牌继续。

下注的威力来自筹码压力。你下注越大，对手继续范围通常越强；你下注越小，对手可以继续的范围通常越宽。但下注大不一定更好，下注小也不一定更弱。尺度必须服务于你的目的。

新手需要先建立一个基本认识：每一次下注都应该是有目的的筹码投入，而不是凭感觉把筹码推出去。

\subsection{下注的两个核心目的：价值与诈唬}

下注的核心目的主要有两个：价值下注和诈唬下注。

价值下注的目标是让更差的牌跟注。你认为自己的牌领先对手相当一部分继续范围，并且对手会用足够多更差牌支付。此时下注的意义是从这些更差牌中获得更多价值。

例如，你持有 \texttt{AK}，翻牌为 \texttt{A72 rainbow}，面对大盲防守。对手范围中有许多较弱 A、小对子和一些浮动牌。你下注时，希望这些更差牌继续。这就是价值下注。

诈唬下注的目标是让更好的牌弃牌。你当前牌力不足以摊牌获胜，或者摊牌价值很低，但你的下注有机会让对手弃掉比你更好的牌。此时下注的意义是通过弃牌权益赢下底池。

例如，你作为翻前加注者，在 \texttt{A72 rainbow} 牌面持有 \texttt{QJ}。你没有成牌，但这个牌面对你的范围更有利。你小注时，可能让大盲的一些小对子、K 高、Q 高或完全空气牌弃掉。这就是诈唬下注。

价值和诈唬不是由你的主观情绪决定的，而是由对手的继续范围决定的。你说“我是在价值下注”，但如果没有更差牌会跟注，那它就不是好的价值下注。你说“我是在诈唬”，但如果没有更好牌会弃牌，那它就不是好的诈唬。

下注前最重要的问题是：

\begin{center}
\textbf{我希望对手跟注，还是希望对手弃牌？}
\end{center}

如果你希望对手跟注，你必须能说出哪些更差牌会跟。
如果你希望对手弃牌，你必须能说出哪些更好牌会弃。
如果你既说不出更差牌会跟，也说不出更好牌会弃，那么这个下注通常没有意义。

\subsection{价值下注的判断标准}

价值下注看似简单，但真正做好并不容易。很多新手拿到强牌时会下注，但他们很少具体思考对手会用什么更差牌跟注。结果有时下注太小，少拿价值；有时下注太大，把更差牌打跑；有时牌力已经不适合价值下注，却仍然盲目继续。

判断一手牌能否价值下注，可以使用以下标准。

第一，你的牌要领先对手足够多的继续范围。不是说你必须永远领先，而是当对手跟注时，你长期应该能赢过足够比例的跟注牌。

第二，对手范围中要有更差牌愿意跟注。价值来自对手错误或边缘支付。如果对手只会用比你强的牌继续，那么你的下注就是在价值切割自己。

第三，下注尺度要能让目标牌继续。薄价值下注通常不适合过大尺度，因为你的目标是让更差牌舒服地支付；强价值牌则可以根据对手承受能力和底池规划使用更大尺度。

第四，你要考虑被加注后的计划。如果你用一手中等牌力下注，面对加注完全不知道怎么办，那么这次下注可能并不舒服。价值下注不是只考虑对手跟注，也要考虑对手反击。

例如，你持有 \texttt{KQ}，公共牌为 \texttt{K8362}，对手一路过牌。你河牌下注小到中等尺度，希望较弱 K、8x、口袋对子支付，这可能是合理薄价值下注。
但如果公共牌为 \texttt{KQJT9}，你只有一对 K，下注后几乎没有更差牌愿意跟注，而大量顺子、两对、强牌会继续，那么下注就不再是好的价值下注。

价值下注的关键不是“我有牌”，而是“更差牌会不会付钱”。

\subsection{薄价值下注}

薄价值下注是价值下注中的重要能力。它指的是你并不持有非常强的牌，但仍然认为对手会用足够多更差牌跟注，因此下注可以长期盈利。

很多新手在河牌价值不足，就是因为他们只敢用坚果或接近坚果的牌下注。只要牌力不是非常强，他们就选择过牌摊牌。这样虽然减少了一些被更强牌跟注的风险，但也错过了大量更差牌支付的机会。

薄价值下注常见于以下情况：

\begin{itemize}
\item 你持有顶对不错踢脚，对手范围中有较弱顶对；
\item 你持有第二强牌，对手会用第三强牌支付；
\item 对手是跟注站，喜欢用边缘牌看摊牌；
\item 河牌没有明显完成强听牌，牌面较安全；
\item 你下注尺度较小，对手容易用更差牌好奇跟注。
\end{itemize}

薄价值下注的核心不是盲目自信，而是精确估算更差跟注牌。如果对手很紧，只用强牌跟注，那么薄价值下注应该减少；如果对手喜欢跟注，那么薄价值下注应该增加。

例如，你在按钮位 open，大盲跟注。公共牌为 \texttt{K8342}，你持有 \texttt{KQ}，对手河牌过牌。面对很多娱乐玩家，你可以下注小到中等尺度，从 \texttt{KJ}、\texttt{KT}、\texttt{K9}、8x 或部分口袋对子中获取价值。
如果你总是过牌摊牌，虽然经常赢，但会少赢一注。

薄价值下注考验的是牌手对对手跟注范围的判断。它不是为了炫技，而是长期盈利的重要来源。

\subsection{诈唬下注的判断标准}

诈唬下注的目标是让更好的牌弃牌。很多新手理解诈唬时，只关注自己“敢不敢打”，却忽略了对手“会不会弃”。真正的诈唬不是胆量问题，而是条件问题。

一个合理诈唬通常需要满足以下条件。

第一，你的牌摊牌价值较低。如果你的牌本来能通过过牌摊牌赢下不少底池，就不一定适合转成诈唬。适合诈唬的牌通常是错过听牌、空气牌或几乎无法摊牌获胜的牌。

第二，你的行动线能够代表强牌。如果你前面行动一直不像强牌，突然河牌大注，对手很容易不相信。好的诈唬需要故事一致。

第三，牌面支持你的范围。如果公共牌明显更有利于对手，你强行诈唬的可信度下降；如果转牌或河牌强化了你的范围，你的诈唬更容易成功。

第四，对手范围中有足够多可以弃掉的更好牌。比如中等对子、弱顶对、错过但仍领先你的 A 高等。如果对手范围要么是强牌，要么是完全空气牌，诈唬价值就会下降。

第五，对手必须有弃牌能力。对跟注站诈唬，是新手最常见的昂贵错误。对不会弃牌的人，你应该少诈唬，多价值下注。

例如你在翻牌和转牌都代表强范围，河牌完成前门同花，而你手中持有一张关键花色 A，阻断对手坚果同花。对手范围中有大量一对和两对类中等牌力。如果对手有弃牌能力，你就可能拥有一个合理的河牌诈唬机会。

反过来，如果对手是“有一对就看到底”的玩家，那么即使牌面再吓人，诈唬也可能没有意义。

诈唬下注的关键不是“我能不能装强”，而是“对手会不会弃掉比我好的牌”。

\subsection{下注尺度如何影响对手继续范围}

下注尺度会直接影响对手的继续范围。一般来说，下注越小，对手可以用越宽的范围继续；下注越大，对手必须用更强、更有韧性的范围继续。

小注常常会被很多边缘牌跟注，例如小对子、弱对、后门听牌、高牌浮动等。它适合用于高频持续下注、薄价值下注、干燥牌面上的范围下注，以及用较低成本给对手施压。

中注通常开始让对手放弃部分弱牌，同时仍然让一些中等牌力和听牌继续。它适合更明确的价值下注、保护性收费和部分半诈唬。

大注会显著压缩对手继续范围。对手通常需要较强成牌、强听牌或较好的抓诈唬牌才能继续。大注更适合极化策略，也就是用强价值牌和高质量诈唬牌下注。

超池下注则会施加更大压力，常见于一方拥有明显坚果优势，或者希望对对手中等牌力施加最大压力的局面。超池下注不是新手必须频繁使用的工具，但需要理解它的含义：它通常代表非常强的价值牌，或者非常明确的诈唬计划。

不同尺度会让对手继续范围发生变化。比如同一手 \texttt{KQ} 顶对，在干燥牌面上小注，可能被许多较弱牌跟注；如果超池下注，对手继续范围会强得多，你的 \texttt{KQ} 可能就不再适合这样下注。

因此，下注尺度不是随便选的，也不是越大越厉害。尺度必须服务于目的：

\begin{itemize}
\item 想让更差牌广泛跟注，可以用较小尺度；
\item 想向听牌和中等牌收费，可以用中等尺度；
\item 想代表极强牌、压迫对手中等牌力，可以用大尺度；
\item 想最大化坚果优势或制造河牌全下压力，可以考虑超池尺度。
\end{itemize}

下注之前，不仅要问“我下注是价值还是诈唬”，还要问“这个尺度是否服务于我的目的”。

\subsection{小注、中注、大注和超池下注的基本含义}

为了让新手更容易理解，可以先把下注尺度粗略分为四类。

\subsubsection*{小注}

小注通常指 1/4 底池到 1/3 底池左右的下注。它的特点是成本低、频率高、对手继续范围宽。

小注适合：

\begin{itemize}
\item 高牌干燥面上的范围持续下注；
\item 薄价值下注；
\item 让对手用大量较弱牌继续；
\item 用较低成本攻击对手空气牌；
\item 控制底池同时保持主动。
\end{itemize}

小注的缺点是压力有限，较难让中等以上牌力弃牌，也不能充分向强听牌收费。

\subsubsection*{中注}

中注通常指 1/2 底池到 2/3 底池左右的下注。它在价值和压力之间较为平衡。

中注适合：

\begin{itemize}
\item 较强价值牌向更差牌收费；
\item 湿润牌面上让听牌付出代价；
\item 半诈唬时同时保留弃牌权益和后续计划；
\item 转牌继续下注时制造适中压力。
\end{itemize}

中注是实战中最常见的尺度之一。它既不会像小注那样压力太低，也不会像大注那样迅速极化范围。

\subsubsection*{大注}

大注通常指 3/4 底池到接近底池的下注。它会明显压缩对手继续范围，通常更适合强价值牌和强诈唬候选。

大注适合：

\begin{itemize}
\item 你拥有坚果优势；
\item 牌面湿润，需要让强听牌支付较高价格；
\item 你希望对对手中等牌力施加压力；
\item 你有强价值牌并希望构建大底池；
\item 你的诈唬牌有阻断效果或较强后续权益。
\end{itemize}

大注的风险是，如果你的牌力不够强，或诈唬条件不足，很容易让更差牌弃掉、更强牌继续。

\subsubsection*{超池下注}

超池下注是指下注金额超过当前底池。它通常表达非常极化的范围：强到希望打巨大底池的价值牌，以及没有摊牌价值、希望最大化弃牌权益的诈唬牌。

超池下注适合：

\begin{itemize}
\item 你拥有明显坚果优势；
\item 对手范围被限制在大量中等牌力；
\item 你可以代表非常强的牌；
\item 你希望对对手一对、两对等牌力施加极大压力；
\item 你有明确河牌计划或全下计划。
\end{itemize}

初学者不需要频繁使用超池下注，但必须明白：超池下注不是情绪化“吓人”，而是基于范围和坚果优势的高压策略。

\subsection{保护性下注为什么不能脱离价值和诈唬}

新手最常说的一句话是：“我下注是为了保护。”

保护这个概念本身并非错误。很多时候，你的牌当前领先，但对手有听牌或高张可以反超你。下注可以让对手为继续支付价格，避免他免费实现权益。

但保护不能成为一个独立于价值和诈唬之外的模糊理由。你必须进一步问：

\begin{itemize}
\item 我保护的是一手什么牌？
\item 对手有哪些更差牌或听牌会跟注？
\item 我的下注是否让这些牌付出了不合适的价格？
\item 如果对手弃牌，我是否成功让有权益的牌放弃？
\item 如果对手加注，我的牌还能不能继续？
\end{itemize}

例如你持有 \texttt{AK}，翻牌为 \texttt{KJ7 two-tone}。你下注可以让较弱 K、Jx、同花听牌、顺子听牌付费。这种保护性下注本质上同时包含价值：更差牌和听牌会继续，你向它们收费。

但如果你持有 \texttt{K9}，同样牌面下注很大，结果更差牌大多弃牌，更强 K、两对、强听牌继续。你以为自己在保护，实际上可能是在用中等牌力把底池做大。

所以，保护必须落回到两个问题：

\begin{center}
\textbf{更差牌是否会跟注？\\
有权益的更好或接近牌是否会弃牌或付出错误价格？}
\end{center}

如果回答不了，所谓保护很可能只是恐惧驱动的下注。

\subsection{半诈唬}

半诈唬是诈唬的一种，但它和纯诈唬不同。纯诈唬通常依赖对手弃牌；半诈唬即使被跟注，也仍然有机会在后续街改善成强牌。

常见半诈唬包括：

\begin{itemize}
\item 同花听牌下注；
\item 两头顺听牌加注；
\item 组合听牌强力进攻；
\item 带后门同花和顺子可能的高牌持续下注；
\item 一对加听牌的转牌加注。
\end{itemize}

半诈唬有两个赢法：

\begin{enumerate}
\item 对手现在弃牌，你直接赢底池；
\item 对手跟注，你在后续街改善成强牌，赢更大底池。
\end{enumerate}

强听牌适合半诈唬，是因为它们在被跟注时仍然保留较高权益。比如坚果同花听牌加两张高牌，不是完全空气牌。它可以通过下注获得弃牌权益，也可以在被跟注后继续改善。

但半诈唬也不能滥用。弱听牌、低同花听牌、没有位置且面对不弃牌对手时，半诈唬价值会下降。对不会弃牌的人，半诈唬更接近单纯追牌；如果价格不合适，就会亏损。

半诈唬的关键是平衡两个因素：

\begin{center}
\textbf{我有没有足够弃牌权益？\\
如果被跟注，我有没有足够牌面权益？}
\end{center}

两者至少要有一个足够强，半诈唬才有意义。

\subsection{试探性下注为什么是错误概念}

娱乐玩家最常见的下注理由之一，是“我下注是为了试探一下”。

试探性下注的问题在于，它把下注目的变得模糊。你下注后，对手跟注、加注或弃牌，确实会给你一些信息。但这些信息并不是免费的，你为它付出了筹码。而且，对手的回应也不一定能准确告诉你答案。

例如你下注，对手跟注。你可能会想：“他跟了，应该有牌。”但他可能是听牌，可能是弱对，可能是慢打强牌，也可能是浮动。
你下注，对手加注。你可能觉得“他很强”，但他也可能在半诈唬。
你下注，对手弃牌。你赢了底池，但你可能本来就领先，也可能把本来会在后续街诈唬的弱牌打走了。

所以，下注不能只是为了“看看他有什么”。正确的思考应该是：

\begin{itemize}
\item 如果我下注，他哪些更差牌会跟？
\item 如果我下注，他哪些更好牌会弃？
\item 如果他加注，我如何应对？
\item 这个下注是否比过牌更盈利？
\end{itemize}

信息是下注的副产品，不是下注的核心目的。你不能为了买一个模糊信息，就随便投入筹码。

\section{新手常见误区}

\subsection{误区一：为了“知道对手有没有”而下注}

这是最典型的试探性下注。新手拿着中等牌力，不确定自己是否领先，于是下注。他们希望通过对手的反应判断对手有没有强牌。

问题是，对手的反应并不总是诚实而清晰。对手跟注不代表一定强，对手弃牌不代表你原本落后，对手加注也不代表一定是坚果。你用下注换来的信息，经常并不准确。

更重要的是，这类下注往往让你的手牌陷入尴尬：更差牌弃牌，更强牌继续，你既没有价值，也没有成功诈唬。

正确做法是：先判断自己的牌力类别。如果是中等摊牌价值，很多时候过牌控池比试探下注更好。

\subsection{误区二：有牌就下注}

“我有牌，所以我要打。”这是新手最自然的想法，但并不总是正确。

你有底对、第二对子、弱顶对，确实是成牌，但它们是否适合下注，要看更差牌是否会跟注。如果下注后只有更强牌继续，那么你的下注就是错误价值下注。

成牌不等于价值下注。价值下注必须有更差跟注范围。

\subsection{误区三：害怕被反超而过度保护}

很多新手在湿润牌面拿到一对后，会立刻大注，理由是“不让他追”。但如果你的牌只是中等牌力，大注可能会产生相反效果：更差牌被你打跑，更强牌和高权益听牌继续。

保护需要考虑牌力和尺度。强牌可以通过下注向听牌收费；中等牌力则要避免因为恐惧而把底池做得过大。

怕被反超不是下注理由。让对手付出错误价格，才是合理保护的基础。

\subsection{误区四：诈唬不看对手类型}

新手有时会迷恋漂亮的诈唬线，却忽略对手是否会弃牌。对跟注站诈唬，是非常低效的策略。对手如果拿着任何对子都想看摊牌，那么你再合理的故事也很难让他弃牌。

面对不同对手，下注目的要调整：

\begin{itemize}
\item 对跟注站，多价值，少诈唬；
\item 对紧弱玩家，可以增加诈唬，减少薄价值；
\item 对激进玩家，可以用中等摊牌价值诱导；
\item 对被动玩家，面对其突然大注要更谨慎。
\end{itemize}

诈唬不是对所有人都一样有效。下注目的必须与对手类型匹配。

\subsection{误区五：下注尺度随意}

有些新手每次下注都用同一个尺度，不管自己是强牌、弱牌、价值还是诈唬。这会让下注无法服务于目标。

如果你想让较弱牌跟注，尺度过大可能把它们赶走。
如果你想让更好牌弃牌，尺度过小可能没有压力。
如果你想向听牌收费，尺度太小会给对手合适价格。
如果你用中等牌力大注，可能只被更强牌继续。

下注尺度必须和目的相匹配。没有目的的尺度，就是随意下注。

\subsection{误区六：半诈唬和乱诈唬混淆}

半诈唬不是随便拿听牌乱打。它需要有牌面权益，也需要有弃牌权益。

如果你的听牌质量很低，对手又不爱弃牌，你的下注就不是高质量半诈唬，而只是用弱牌扩大底池。反过来，如果你有强听牌、范围优势和位置，半诈唬可以成为非常有力的进攻方式。

半诈唬的关键是：被跟注后仍然有机会，下注时也有机会让对手弃牌。

\section{实战思考框架}

本章的核心框架是“下注前必问”。

\begin{center}
\textbf{如果我是价值下注，哪些更差牌会跟？\\
如果我是诈唬，哪些更好牌会弃？\\
如果两个问题都回答不了，就不要轻易下注。}
\end{center}

\subsection{第一问：如果我是价值下注，哪些更差牌会跟？}

当你准备下注时，首先判断自己是否在价值下注。

如果是价值下注，你必须具体列出更差跟注牌。例如：

\begin{itemize}
\item 我有强顶对，对手会用较弱顶对跟注；
\item 我有两对，对手会用顶对或弱两对跟注；
\item 我有坚果同花，对手会用较小同花或暗三条跟注；
\item 我有薄价值牌，对手会用第二对子或抓诈唬牌跟注。
\end{itemize}

如果你只能说“我可能领先”，但说不出哪些更差牌会跟注，那么下注理由不够充分。

价值下注还要考虑尺度。你要让目标更差牌愿意支付。如果目标牌很弱，你的尺度应该更容易被跟注；如果目标牌较强，或者对手是跟注站，你可以使用更大尺度。

\subsection{第二问：如果我是诈唬，哪些更好牌会弃？}

如果你的下注不是价值下注，那就要判断是否是诈唬。

诈唬下注必须能让更好牌弃牌。你需要具体列出目标弃牌范围。例如：

\begin{itemize}
\item 我希望对手弃掉小对子；
\item 我希望对手弃掉弱顶对；
\item 我希望对手弃掉错过但仍领先我的 A 高；
\item 我希望对手弃掉中等摊牌价值牌。
\end{itemize}

同时，你要判断对手是否真的会弃这些牌。一个理论上应该弃牌但实际上从不弃牌的对手，不是好的诈唬对象。

诈唬还要考虑你的行动线是否可信。你能否代表完成的同花、顺子、两对或强顶对？如果你的故事不一致，对手更容易抓你。

\subsection{第三问：如果两个问题都回答不了，就不要轻易下注}

很多糟糕下注既不是价值，也不是诈唬。

例如你持有中等对子，下注后更差牌大多弃牌，更强牌不会弃牌。这个下注没有价值，也没有诈唬效果。它只是让你投入筹码，却没有让对手犯错。

面对这种局面，过牌通常更好。过牌可以控制底池，保留摊牌价值，观察对手行动，避免把自己放进尴尬位置。

新手要建立一个纪律：

\begin{center}
\textbf{下注之前说不清目的，就先不要下注。}
\end{center}

扑克中不下注也是一种主动选择。过牌不是软弱，弃牌也不是失败。没有目的的下注，才是真正危险的行动。

\subsection{下注目的与牌力类型的对应关系}

不同牌力类型，对应不同下注目的。

\begin{center}
\begin{tabular}{lll}
\toprule
\textbf{牌力类型} & \textbf{常见下注目的} & \textbf{注意事项} \\
\midrule
坚果牌 & 强价值下注 & 构建大底池，避免过度慢打 \\
强价值牌 & 多街价值下注 & 确认更差牌会继续 \\
薄价值牌 & 小到中等价值下注 & 尺度要让更差牌愿意跟 \\
中等摊牌价值 & 多数不宜大注 & 控池，避免价值切割自己 \\
强听牌 & 半诈唬 & 同时依赖弃牌权益和牌面权益 \\
弱听牌 & 谨慎继续 & 不宜用大注硬打 \\
空气牌 & 选择性诈唬 & 必须有弃牌对象和可信故事 \\
\bottomrule
\end{tabular}
\end{center}

这张表可以帮助你把第六章的牌力功能，与本章的下注目的结合起来。

\subsection{下注尺度选择的简化流程}

选择下注尺度时，可以按以下流程思考：

\begin{enumerate}
\item 我是价值还是诈唬？
\item 我的目标牌是强、中等还是弱？
\item 我希望对手继续范围宽一点，还是窄一点？
\item 牌面是干燥还是湿润？
\item 这次下注是否要为下一条街构建底池？
\end{enumerate}

如果你是薄价值，希望较弱牌跟注，通常使用较小尺度。
如果你是强价值，希望构建大底池，可以使用中大尺度。
如果你是半诈唬，希望兼顾弃牌权益和牌面权益，可以根据牌面选择中等或较大尺度。
如果你是纯诈唬，希望让中等牌力弃牌，尺度需要足够有压力。
如果你不知道尺度为什么这样选，说明你还没有明确下注目的。

\section{牌例拆解}

\subsection{牌例一：强价值下注}

\subsubsection*{牌局背景}

有效筹码 100BB。你在按钮位拿到 \texttt{AKo}，前面玩家弃牌，你 open raise 到 2.5BB。大盲跟注。

翻牌为 \texttt{A72 rainbow}。大盲过牌，行动轮到你。

\subsubsection*{新手常见想法}

新手可能会想：

\begin{itemize}
\item “我中了顶对，下注。”
\item “牌面安全，可以打。”
\item “不能让他免费看牌。”
\end{itemize}

这些想法方向没错，但需要更明确下注目的。

\subsubsection*{理性分析}

第一，判断牌力。你持有 \texttt{AK}，在 \texttt{A72 rainbow} 上是强顶对好踢脚，属于强价值牌。

第二，判断下注目的。这是价值下注。你希望大盲用较弱 A、7x、小口袋对子、部分高牌浮动继续。

第三，判断牌面结构。牌面干燥，听牌很少。你不需要为了保护而使用大尺度。小注即可让大量更差牌继续，同时保持范围优势。

第四，选择尺度。1/3 底池左右的小注通常合理。它能让更差牌跟注，也能让对手大量空气牌弃牌。

\subsubsection*{推荐行动}

下注小尺度，进行价值下注。

\subsubsection*{本手牌教训}

强价值下注不是因为“我有牌所以打”，而是因为对手有足够多更差牌会跟注。尺度应服务于让这些牌继续。

\subsection{牌例二：中等牌力不适合试探下注}

\subsubsection*{牌局背景}

有效筹码 100BB。你在 cutoff 拿到 \texttt{JTs}，open raise 到 2.5BB。大盲跟注。

翻牌为 \texttt{KJ7 two-tone}。大盲过牌，行动轮到你。

\subsubsection*{新手常见想法}

很多新手会想：

\begin{itemize}
\item “我有第二对子，下注看看他有没有 K。”
\item “牌面有听牌，不能给他免费看。”
\item “如果他跟注，我再判断。”
\end{itemize}

这就是典型的试探性下注。

\subsubsection*{理性分析}

第一，判断牌力。你持有第二对子，有一定摊牌价值，但不是强价值牌。

第二，判断价值下注。下注后，哪些更差牌会跟注？可能有一些 7x、听牌或较弱 Jx，但很多更差空气牌会弃牌。更强 Kx、强 J、两对、暗三条和强听牌会继续。

第三，判断诈唬下注。你下注能让哪些更好牌弃牌？大多数 Kx 不会因为一枪就弃牌。更强 J 也未必弃。你的下注很难让足够多更好牌弃牌。

第四，比较过牌。过牌可以控制底池，保留摊牌价值，也让对手的空气牌和错过听牌留在范围里。

\subsubsection*{推荐行动}

倾向过牌控池。

如果对手非常被动且会用大量更差牌跟注，可以考虑小注；但默认情况下，不应为了“试探”而下注。

\subsubsection*{本手牌教训}

中等摊牌价值最容易被试探性下注误用。下注前必须问：更差牌会跟吗？更好牌会弃吗？如果都不明确，过牌通常更好。

\subsection{牌例三：半诈唬下注}

\subsubsection*{牌局背景}

有效筹码 100BB。你在按钮位拿到 \texttt{AQs}，前面玩家弃牌，你 open raise 到 2.5BB。大盲跟注。

翻牌为 \texttt{J T 4}，其中两张与你同花。大盲过牌，行动轮到你。

\subsubsection*{新手常见想法}

一些新手会想：

\begin{itemize}
\item “我还没中牌，等中了再打。”
\item “如果下注被跟，我现在还是 A 高。”
\item “听牌就跟，不要主动打。”
\end{itemize}

这些想法低估了半诈唬的价值。

\subsubsection*{理性分析}

第一，判断牌力。你有坚果同花听牌、两张高牌和顺子相关可能，是强听牌。

第二，判断下注目的。这不是纯价值下注，而是半诈唬。你希望对手弃掉部分小对子、弱 4x、没有强权益的空气牌；如果对手跟注，你仍然有大量转牌和河牌改善机会。

第三，判断牌面。牌面较湿润，对手有继续范围，但你的牌权益很高，适合主动进攻。

第四，选择尺度。可以使用中等尺度，既制造弃牌权益，也为后续完成强牌构建底池。

\subsubsection*{推荐行动}

下注半诈唬。

\subsubsection*{本手牌教训}

半诈唬不是乱 bluff。强听牌下注有两个赢法：对手现在弃牌，或者你后续改善成强牌。

\subsection{牌例四：薄价值下注}

\subsubsection*{牌局背景}

有效筹码 100BB。你在按钮位拿到 \texttt{KQ}，open raise 到 2.5BB。大盲跟注。

翻牌为 \texttt{K83 rainbow}。大盲过牌，你下注 1/3 底池，大盲跟注。

转牌为 \texttt{6}。双方过牌。

河牌为 \texttt{2}。大盲过牌，行动轮到你。

\subsubsection*{新手常见想法}

新手可能会想：

\begin{itemize}
\item “我有顶对，过牌也能赢。”
\item “如果他有两对或暗三条，我下注会输钱。”
\item “河牌没必要冒险。”
\end{itemize}

这类想法会导致漏掉薄价值。

\subsubsection*{理性分析}

第一，判断牌力。你持有 \texttt{KQ}，在 \texttt{K8362} 上是较强顶对。转牌双方过牌后，对手范围并不一定很强。

第二，判断价值下注。更差牌包括较弱 K、8x、部分口袋对子，甚至一些好奇跟注的 A 高。面对小到中等尺度，这些牌可能会支付。

第三，判断尺度。因为这是薄价值下注，尺度不宜过大。小到中等下注更容易被更差牌跟注。

第四，判断风险。你会输给少量两对和暗三条，但这不代表不能下注。价值下注不是要求永远领先，而是要求被跟注时长期领先足够多。

\subsubsection*{推荐行动}

下注小到中等尺度，进行薄价值下注。

\subsubsection*{本手牌教训}

过牌能赢，不代表过牌最好。如果更差牌会支付，河牌应该争取薄价值。

\subsection{牌例五：错误的河牌诈唬对象}

\subsubsection*{牌局背景}

有效筹码 100BB。你在按钮位拿到 \texttt{QTs}，open raise 到 2.5BB。大盲是一名明显跟注站，喜欢用任何对子跟到河牌。

翻牌为 \texttt{K72 rainbow}。大盲过牌，你下注 1/3 底池，大盲跟注。

转牌为 \texttt{4}。大盲过牌，你继续下注 2/3 底池，大盲跟注。

河牌为 \texttt{9}。你没有成牌。大盲过牌，行动轮到你。

\subsubsection*{新手常见想法}

新手可能会想：

\begin{itemize}
\item “我已经打了两枪，第三枪才能把故事讲完。”
\item “他可能只有 7 或小对子，河牌大注能打走。”
\item “我不打就一定输了。”
\end{itemize}

这些想法忽略了对手类型。

\subsubsection*{理性分析}

第一，判断摊牌价值。你没有成牌，几乎没有摊牌价值。

第二，判断诈唬目标。你希望对手弃掉 7x、小对子、弱 K 或其他比你好的牌。

第三，判断对手是否会弃。关键问题是：这个大盲是跟注站。他已经用很多边缘牌跟到河牌，且倾向于用任何对子看摊牌。对这样的对手，诈唬成功率很低。

第四，判断行动线。虽然你可以代表部分 Kx 和强牌，但对手如果不关心你的故事，诈唬就没有意义。

\subsubsection*{推荐行动}

过牌放弃。

\subsubsection*{本手牌教训}

诈唬必须选择会弃牌的对象。对不会弃牌的人，少诈唬，多价值下注。好的故事讲给听得懂的人才有意义。

\section{本章总结}

本章的核心，是让你真正理解下注不是凭感觉把筹码推出去，而是带有明确目的的进攻行动。

筹码是你唯一的武器。你通过下注表达范围、改变价格、迫使对手筛选继续范围。但筹码不能随意使用，每一次下注都应该服务于具体目标。

下注的两个核心目的，是价值和诈唬。

价值下注的目标，是让更差牌跟注。
诈唬下注的目标，是让更好牌弃牌。

保护、半诈唬、薄价值、施压、阻止免费看牌等概念，都不能脱离这两个核心目的独立存在。保护本质上要么是让更差牌和听牌付费，要么是迫使有权益的牌放弃；半诈唬是带有牌面权益的诈唬；薄价值是边缘但仍然盈利的价值下注。

下注尺度必须服务于下注目的。小注让对手继续范围更宽，适合薄价值、高频下注和低成本施压；中注平衡价值与压力；大注压缩对手继续范围，更适合强价值和极化诈唬；超池下注则通常需要坚果优势和明确的极化策略。

本章的核心框架是：

\begin{center}
\textbf{如果我是价值下注，哪些更差牌会跟？\\
如果我是诈唬，哪些更好牌会弃？\\
如果两个问题都回答不了，就不要轻易下注。}
\end{center}

新手最需要戒掉的是试探性下注。下注不是为了“看看对手有没有”，因为对手的反应并不总是清晰可靠。信息只是下注的副产品，不是下注的目的。

下一章我们将讨论“激进的游戏”。当你理解下注目的之后，就要进一步理解主动权、弃牌权益和进攻频率。德州扑克不是只靠摊牌赢底池的游戏，优秀牌手必须学会在合适条件下主动争取底池。

\section{课后训练}

\subsection*{训练一：下注目的记录}

下一次打牌后，记录 10 手你主动下注的牌。每一手都回答：

\begin{enumerate}
\item 我下注是价值还是诈唬？
\item 如果是价值下注，哪些更差牌会跟注？
\item 如果是诈唬下注，哪些更好牌会弃牌？
\item 我的下注尺度是否服务于这个目的？
\item 如果对手加注，我当时有没有计划？
\end{enumerate}

如果你无法回答第 2 或第 3 个问题，说明这次下注目的不清。

\subsection*{训练二：找出试探性下注}

从最近牌局中找出 5 手你“为了看看对手有没有”而下注的牌。

逐手分析：

\begin{itemize}
\item 我的牌属于什么牌力类型？
\item 更差牌是否会跟注？
\item 更好牌是否会弃牌？
\item 过牌是否更合理？
\item 下注后我获得的信息是否真的可靠？
\end{itemize}

这个训练的目标，是逐步戒掉试探性下注。

\subsection*{训练三：薄价值下注训练}

找出 10 手你河牌过牌摊牌并赢下底池的牌，判断其中哪些可以薄价值下注。

重点回答：

\begin{enumerate}
\item 对手是否可能用更差牌跟注？
\item 我应该使用小注、中注还是大注？
\item 对手是跟注站、紧弱玩家，还是正常玩家？
\item 我当时过牌是合理控池，还是害怕下注？
\end{enumerate}

训练目标是提高河牌价值获取能力。

\subsection*{训练四：诈唬对象选择}

记录 5 手你尝试诈唬的牌，逐手回答：

\begin{itemize}
\item 对手是否有弃牌能力？
\item 我的行动线是否能代表强牌？
\item 牌面是否支持我的范围？
\item 我的牌是否几乎没有摊牌价值？
\item 我的下注尺度是否足够让更好牌弃牌？
\end{itemize}

如果对手不具备弃牌能力，未来应减少类似诈唬。

\subsection*{训练五：下注尺度复盘}

找出最近 10 次下注，记录你的下注尺度，并判断它是否符合目的：

\begin{itemize}
\item 小注：是否为了让更宽范围继续或低成本施压？
\item 中注：是否为了价值和保护之间平衡？
\item 大注：是否因为我有强价值或强诈唬候选？
\item 超池下注：是否真的有坚果优势和极化理由？
\end{itemize}

这个训练的目标，是让下注尺度从习惯动作变成策略选择。

\subsection*{训练六：半诈唬判断}

记录 10 手你持有听牌并主动下注或加注的牌，分析它们是否是真正的半诈唬。

逐手回答：

\begin{enumerate}
\item 我是否有足够牌面权益？
\item 对手是否有足够弃牌率？
\item 我的听牌是强听牌还是弱听牌？
\item 如果被跟注，转牌或河牌有哪些好牌？
\item 如果对手加注，我是否有继续计划？
\end{enumerate}

不要把所有听牌下注都称为半诈唬。只有同时具备弃牌权益和牌面权益的下注，才是高质量半诈唬。

\chapter{激进的游戏}

\section{本章导入}

德州扑克不是被动等待好牌的游戏。很多新手刚开始打牌时，最自然的策略是等待强牌：拿到好牌就下注，没牌就过牌或弃牌；成牌就继续，没中就放弃。这样的打法表面上稳健，实际上会让牌手失去一半赢底池的方式。

获得底池通常有两种方法：

\begin{enumerate}
\item 摊牌时拥有最好牌；
\item 通过进攻让对手弃牌。
\end{enumerate}

新手往往只会第一种方式。他们必须等自己成牌，必须等自己领先，必须等公共牌帮助自己，才敢投入筹码。这样一来，很多原本可以通过主动进攻赢下的底池，都被他们放弃了。更重要的是，对手会很容易读懂他们：下注就是有牌，过牌就是放弃。这样的玩家很难给对手制造持续压力。

优秀牌手不只是等好牌，而是在合适的局面主动争取底池。他们理解主动权，理解弃牌权益，理解什么牌适合半诈唬，理解什么牌面适合高频进攻，也理解什么对手可以被施压。真正的激进不是乱 bluff，而是有条件、有目标、有计划地施压。

激进并不等于鲁莽。鲁莽是没有范围优势也进攻，没有牌面支持也进攻，对手不会弃牌也进攻，自己的牌没有功能也进攻。有效激进则不同，它建立在三个条件之上：

\begin{center}
\textbf{我的范围有优势。\\
我的牌有合适的功能。\\
对手有足够多可以弃掉的牌。}
\end{center}

这就是本章的核心框架。

本章要解决的问题包括：为什么不能只等好牌？什么是弃牌权益？激进和鲁莽有什么区别？什么情况下应该主动抢夺底池？什么情况下激进反而是亏损？加注除了代表强牌之外，还有哪些意义？

学完本章后，你应该逐渐从“只有牌才敢打”的被动玩家，转向“在合适条件下主动争取底池”的思考型玩家。

\section{核心概念}

\subsection{为什么扑克需要激进性}

德州扑克中，如果你只靠摊牌赢底池，就会天然处于劣势。原因很简单：多数时候，你并不会在翻牌、转牌或河牌都拿到明显强牌。很多手牌会错过牌面，很多成牌只是中等牌力，很多听牌最终不会完成。如果你只能在自己牌力很强时才下注，那么你的赢池方式会非常有限。

激进性的重要意义在于，它让你拥有第二种赢法：让对手弃牌。

例如，你在按钮位 open，大盲跟注。翻牌为 \hand{A♠ 7♣ 2♦ rainbow}。你手持 \hand{Q♠ J♣}，没有成牌。如果你只看自己的牌，你会选择过牌放弃。但这个牌面对你的按钮 open range 更有利，大盲有许多没有击中的牌。你的小尺度持续下注可能直接赢下底池。此时你并不是靠摊牌赢，而是靠范围优势和弃牌权益赢。

如果一个玩家完全缺乏激进性，对手可以轻松应对他：

\begin{itemize}
\item 他下注时，对手知道他通常有牌；
\item 他过牌时，对手可以频繁下注抢夺；
\item 他拿中等牌力时，经常被迫猜测；
\item 他错过牌面时，几乎自动放弃；
\item 他很难从对手的弃牌中获利。
\end{itemize}

扑克中的盈利来自两个方向：一是强牌拿价值，二是让对手在不舒服的范围中弃牌。没有激进性，就相当于只用一条腿走路。

但激进性必须建立在理性基础上。有效激进不是更多下注，而是更准确地下注；不是更多 bluff，而是选择更好的 bluff；不是永远进攻，而是在条件有利时进攻。

\subsection{主动权}

主动权通常来自翻前或前一街的加注、下注行动。拥有主动权的一方，在后续街更容易代表强范围，也更容易通过持续下注给对手施压。

例如你翻前 open raise，大盲跟注。你是翻前进攻方，翻后许多牌面上你可以继续下注。对手作为防守方，必须面对你的下注决定是否继续。这种结构本身就给了你主动权。

主动权的价值主要体现在三个方面。

第一，它让你拥有先发压力。你可以通过下注迫使对手筛选范围，而不是被动等待对手下注。

第二，它让你的范围故事更连贯。翻前加注、翻牌下注、转牌继续，这些行动可以共同表达强范围。如果牌面支持你的范围，对手的中等牌力会承受越来越大压力。

第三，它让你拥有更多赢池方式。你可以摊牌赢，也可以通过下注让对手弃牌；而被动玩家更多只能依赖摊牌。

不过，主动权不是无限开火的许可证。你是翻前加注者，不代表每个翻牌都必须 C-bet；你翻牌下注了，不代表每个转牌都必须开第二枪。主动权只是给你进攻资格，是否进攻仍然要看范围优势、牌面结构、手牌功能和对手类型。

\subsection{弃牌权益}

弃牌权益，指的是你通过下注或加注让对手弃牌，从而直接赢下底池的可能性。它是激进策略的核心。

如果你下注后，对手有一定概率弃牌，那么即使你的牌当前不是最好牌，你仍然可以通过下注盈利。尤其当你的牌同时拥有一定牌面权益时，下注会变得更有价值。

例如你持有同花听牌。你下注后，对手可能弃牌，你直接赢底池；如果对手跟注，你仍然有机会在后续街完成同花。这就是弃牌权益和牌面权益叠加带来的力量。

弃牌权益取决于几个因素：

\begin{itemize}
\item 对手范围中有多少弱牌和中等牌力；
\item 你的下注尺度是否足够制造压力；
\item 你的行动线是否可信；
\item 牌面是否支持你代表强牌；
\item 对手是否有弃牌能力。
\end{itemize}

新手经常忽略弃牌权益，所以他们会觉得“我现在没有牌，下注就是乱打”。但事实上，很多下注并不是因为你当前已经领先，而是因为你可以让对手放弃更好牌，或者放弃仍有权益的牌。

不过，弃牌权益不是幻想出来的。对不会弃牌的玩家，弃牌权益很低。面对跟注站，诈唬和半诈唬价值下降，价值下注价值上升。激进之前，必须判断对手是否真的会弃牌。

\subsection{激进和鲁莽的区别}

激进和鲁莽表面上都表现为下注、加注、施压，但本质完全不同。

激进是有条件的进攻。你进攻之前，知道自己为什么下注，知道自己代表什么范围，知道对手有哪些牌会弃，知道如果被跟注或加注应该怎么办。激进的目标是让对手犯错：用更差牌付钱，或者弃掉本该继续的权益。

鲁莽是没有条件的进攻。牌面不支持自己也打，对手不弃牌也打，手牌没有权益也打，被加注后没有计划也打。鲁莽的目标往往不是长期 EV，而是情绪释放、证明自己、追回损失或不想显得软弱。

可以用以下表格区分两者：

\begin{center}
\begin{tabular}{lll}
\toprule
\textbf{比较维度} & \textbf{有效激进} & \textbf{鲁莽进攻} \\
\midrule
范围基础 & 有范围优势或坚果优势 & 不考虑范围 \\
手牌功能 & 价值牌、强听牌、好诈唬候选 & 任意牌硬打 \\
对手类型 & 对手有可弃牌范围 & 对跟注站强行 bluff \\
下注目的 & 价值或诈唬明确 & 试探、发泄、追损 \\
后续计划 & 知道转牌、河牌如何继续 & 被跟注后大脑空白 \\
\bottomrule
\end{tabular}
\end{center}

所以，本章强调激进，并不是鼓励新手乱打。恰恰相反，真正的激进比被动更需要纪律。因为你投入筹码更多，错误成本也更高。如果没有条件和计划，所谓激进只会变成更快输钱。

\subsection{进攻不是乱 bluff，而是有条件地施压}

很多新手一听到“激进”，就以为是多诈唬、多开火、多加注。其实进攻并不等于乱 bluff。进攻的本质是施压，而施压必须建立在对手承受压力的能力有限这一事实之上。

当对手范围中有大量中等牌力时，进攻最有效。因为中等牌力最怕面对多街压力。它们能击败一些弱牌和诈唬，但无法舒服地承受大底池。如果你的下注线能够代表强牌，对手就会被迫弃掉一部分仍然有摊牌价值的牌。

例如，对手在大盲跟注你的按钮 open。翻牌 \hand{A♠ 7♣ 2♦ rainbow}，他 check-call 小注。转牌 \hand{K}，他再次过牌。这个 K 强化你的范围。你继续下注，可以让他的 7x、小对子、弱 A 甚至部分 Kx 承受压力。这里的进攻不是乱 bluff，而是利用范围优势和转牌优势施压。

但如果对手范围非常强，或者牌面明显击中对手，强行进攻就会变差。例如你在前位 open，大盲跟注，翻牌 \hand{9♠ 8♠ 7♦ two-tone}，转牌 \hand{6}。这类牌面极大强化大盲范围。你如果拿空气牌继续大额下注，往往不是有效激进，而是鲁莽。

进攻前要问：

\begin{center}
\textbf{对手有哪些牌会因为我的下注而难受？}
\end{center}

如果对手大多是强牌，不会难受；如果对手是跟注站，不愿弃牌；如果你的故事不可信，那么进攻价值就很低。

\subsection{什么牌适合半诈唬}

半诈唬是激进游戏中最重要的工具之一。它介于价值下注和纯诈唬之间：你当前未必领先，但如果对手弃牌，你直接赢；如果对手跟注，你仍然有机会改善成强牌。

适合半诈唬的牌通常具备以下特点：

\begin{itemize}
\item 有较强牌面权益，例如坚果同花听牌、两头顺听牌、组合听牌；
\item 有额外高张价值，例如同花听牌加两个高张；
\item 有后门权益，例如后门同花加后门顺子；
\item 有阻断效果，例如持有关键 A 或 K，减少对手强牌组合；
\item 被跟注后仍然有后续计划。
\end{itemize}

例如，你持有 \hand{A♠ Q♠}，翻牌为 \hand{J♠ T♣ 4♦}，其中两张与你同花。这是非常好的半诈唬候选。你有坚果同花听牌、两张高牌以及顺子潜力。下注后，对手可能弃掉一些小对子或弱牌；即使被跟注，你也有许多好转牌。

相反，弱卡顺、低同花听牌、没有额外权益的低质量听牌，并不适合频繁大规模半诈唬。它们被跟注后权益不足，完成后也可能被更强牌压制。

半诈唬不是“有听牌就打”，而是“有足够牌面权益和弃牌权益时打”。

\subsection{什么牌面适合高频进攻}

并不是所有牌面都适合高频进攻。牌面是否适合进攻，取决于它对谁的范围更有利，以及它是否容易让对手范围错过。

适合高频进攻的常见牌面包括：

\begin{itemize}
\item 高牌干燥面，例如 \hand{A♠ 7♣ 2♦ rainbow}、\hand{K♠ 8♣ 3♦ rainbow}；
\item 明显有利于翻前加注者的牌面；
\item 对手防守范围中有大量没有击中的牌；
\item 听牌较少、对手难以强力反击的牌面；
\item 转牌或河牌强化你范围的牌面。
\end{itemize}

例如你翻前 open，大盲跟注。翻牌 \hand{A♠ 7♣ 2♦ rainbow}，这个牌面对你的 open range 通常更有利。你可以用较宽范围小注进攻，包括强 A、弱 A、高对、部分空气牌和后门权益牌。

不适合盲目高频进攻的牌面包括：

\begin{itemize}
\item 中低连接面，例如 \hand{9♠ 8♣ 7♦}、\hand{8♠ 7♣ 6♦}、\hand{T♠ 9♣ 8♦}；
\item 两张同花且顺子听牌丰富的湿润面；
\item 更容易击中大盲防守范围的牌面；
\item 对手拥有更多两对、暗三条和顺子组合的牌面；
\item 多人底池中的复杂牌面。
\end{itemize}

在这些牌面上，进攻方需要更谨慎。即使你是翻前加注者，也不能自动持续下注。

\subsection{什么对手适合被施压}

激进策略必须考虑对手类型。不是所有对手都适合被施压，也不是所有对手都应该被频繁诈唬。

适合被施压的对手通常有以下特点：

\begin{itemize}
\item 翻前入池较宽，翻后范围较弱；
\item 翻牌喜欢跟注，但转牌和河牌面对压力会弃牌；
\item 害怕大底池；
\item 中等牌力不愿意承受多街下注；
\item 能够理解牌面威胁，并愿意弃牌。
\end{itemize}

这类对手经常在翻牌跟得较宽，但转牌或河牌面对大注会弃掉大量中等牌力。对他们，合理的第二枪和第三枪会有较好效果。

不适合频繁施压的对手包括：

\begin{itemize}
\item 跟注站：有任何对子都想看摊牌；
\item 极端激进玩家：频繁反加，喜欢对抗；
\item 范围很紧的玩家：继续到后面街时通常很强；
\item 不理解牌面威胁的娱乐玩家：不会因为牌面变危险而弃牌。
\end{itemize}

面对跟注站，最好的策略不是复杂诈唬，而是用价值牌多下注。面对过紧玩家，可以用更高频率偷池和施压。面对激进玩家，则要选择更强、更能承受压力的范围继续。

有效激进不是固定风格，而是针对对手漏洞的调整。

\subsection{什么情况下不该盲目激进}

激进有价值，但盲目激进会迅速亏损。以下情况通常不适合强行进攻。

第一，对手不会弃牌。无论你的故事多合理，对手如果拿一对就跟到底，诈唬价值都会大幅下降。

第二，牌面明显有利于对手范围。例如低连接湿润牌面、完成对手主要听牌的转牌、多人底池中的协调牌面。

第三，你的牌没有功能。没有摊牌价值、没有听牌、没有阻断效果、没有后门权益的纯空气牌，在不利牌面上强行下注，通常质量很低。

第四，你没有后续计划。翻牌下注被跟注后，转牌不知道怎么办；转牌下注被跟注后，河牌不知道是否继续。没有计划的进攻容易变成孤立下注。

第五，你是情绪驱动。刚被对手诈唬、刚输大底池、想追回损失、想证明自己，这些都不是进攻理由。

第六，多人底池。多人底池中，对手总范围更容易有人击中强牌，弃牌权益下降，诈唬成功率降低。新手在多人底池中应该显著减少无牌进攻。

激进之前，一定要确认自己不是在用“进攻”掩盖情绪。

\subsection{加注的意义}

加注是比下注更强的进攻行动。它不仅投入更多筹码，也会进一步压缩对手范围。很多新手把加注理解为“我有强牌”，但实际上，加注可以有多种意义。

第一，价值加注。你持有强牌，希望更差牌继续支付。例如你在翻牌中暗三条，面对对手下注，可以加注构建底池。

第二，保护性加注。当你当前领先但牌面湿润，对手可能有许多听牌时，加注可以让对手为继续支付更高价格。不过它仍然不能脱离价值逻辑：你要确认更差牌和听牌会继续。

第三，半诈唬加注。你持有强听牌，通过加注获得弃牌权益。如果对手弃牌，你直接赢；如果对手跟注，你仍然有改善机会。

第四，施压加注。有时你判断对手下注范围较宽，包含大量中等牌力或空气牌，通过加注可以迫使他弃掉一部分范围。

第五，隔离和争取主动权。翻前 3-bet、翻牌加注，都可能让你从被动防守转为主动进攻。

但加注也有风险。它会迅速放大底池，并让对手继续范围变强。如果你用中等牌力随意加注，往往会让更差牌弃牌、更强牌继续。新手使用加注时，必须先明确：我是价值、半诈唬，还是施压？如果说不清，就不要随便加注。

\section{新手常见误区}

\subsection{误区一：只等好牌}

很多新手过于被动，认为只有拿到强牌才能赢。他们错过牌面就放弃，中等牌力只想便宜摊牌，听牌只会跟注等待完成。

这种打法的问题是，赢池方式太单一。你必须依赖发牌帮助，必须摊牌时最好，才能赢底池。对手一旦发现你只在有牌时下注，就可以轻松弃掉弱牌，或者在你过牌时频繁抢池。

德州扑克需要主动争取底池。你不能只等待自己拥有最好牌，也要学会在范围、牌面和对手都合适时，用进攻让对手弃牌。

\subsection{误区二：把激进理解为乱 bluff}

另一类新手走向相反极端。他们听说“要激进”，就开始频繁 bluff，任何牌都想打走对手。结果对手一跟注，他们就没有计划；对手一加注，他们就陷入混乱。

激进不是乱 bluff。激进需要满足条件：范围有优势，牌有功能，对手有可弃牌范围。缺少这些条件，所谓激进只是鲁莽。

真正的激进不是为了显得强，而是为了让对手在错误范围中做困难决策。

\subsection{误区三：该进攻时被动}

有些新手拿到强听牌，却只是被动跟注；处在有利牌面，却因为自己没有成牌而过牌；面对容易弃牌的对手，却不敢开第二枪。

这类错误看似保守，实际上会损失大量 EV。强听牌通过半诈唬可以拥有两个赢法，而被动跟注通常只剩下完成听牌这一条路。有范围优势的牌面，如果你不进攻，就放弃了对对手空气牌和中等牌力施压的机会。

该进攻时被动，会让你变成只能靠摊牌赢的人。

\subsection{误区四：该放弃时硬冲}

也有新手在完全不合适的局面硬冲。牌面明显击中对手，对手是跟注站，自己没有任何权益，却仍然连续下注。这样的打法不是勇敢，而是亏损。

扑克不是每个底池都要争。优秀牌手知道哪些底池值得抢，也知道哪些底池应该放弃。放弃低质量进攻机会，才能保留筹码在高质量机会中施压。

\subsection{误区五：忽略对手类型}

有些玩家使用同一套进攻策略对所有人。对紧弱玩家不敢施压，对跟注站拼命 bluff，对激进玩家用弱牌硬抗。这些都是没有根据对手类型调整的表现。

对手不同，弃牌权益不同。
弃牌权益不同，激进策略就不同。

面对会弃牌的人，进攻可以增加；面对不会弃牌的人，价值下注才是核心。激进不是固定风格，而是针对对手漏洞的策略。

\subsection{误区六：加注没有目的}

新手有时会用加注“看看对手是不是真有牌”，或者因为“不想被欺负”而加注。这和第七章中的试探性下注一样，都是目的不清。

加注必须更谨慎，因为它比普通下注投入更多筹码。加注前要问：

\begin{itemize}
\item 我是价值加注吗？
\item 我是半诈唬加注吗？
\item 我希望对手弃掉哪些牌？
\item 如果对手再加注，我怎么办？
\item 我的牌是否适合打大底池？
\end{itemize}

如果这些问题回答不了，就不要用加注表达情绪。

\section{实战思考框架}

本章的核心框架是“有效激进三条件”：

\begin{center}
\textbf{我的范围有优势。\\
我的牌有合适的功能。\\
对手有足够多可以弃掉的牌。}
\end{center}

\subsection{第一条件：我的范围有优势}

有效进攻首先要看范围优势。你不是因为“想打”而打，而是因为这个牌面、这条行动线或这张转牌更支持你的范围。

例如你翻前 open，大盲跟注，翻牌为 \hand{A♠ 7♣ 2♦ rainbow}。这个牌面通常更有利于翻前加注者，因为你有更多强 A、大对子和高牌组合。你可以更高频率进攻。

再如你翻牌下注被跟注，转牌出现 \hand{K}，而这张 K 明显强化你的翻前强范围。你可以用强价值牌继续，也可以选择合适诈唬牌开第二枪。

相反，如果牌面是 \hand{9♠ 8♠ 7♦ two-tone}、\hand{8♠ 7♣ 6♦}、\hand{T♠ 9♣ 8♦} 这类更容易击中防守方范围的结构，你的范围优势下降，进攻频率应降低。

判断范围优势时，可以问：

\begin{itemize}
\item 这个牌面更容易击中谁的翻前范围？
\item 谁有更多顶对、超对、暗三条、两对、顺子？
\item 这张转牌或河牌强化了谁的故事？
\item 我的前面行动是否能自然代表强牌？
\end{itemize}

没有范围基础的进攻，很容易被对手继续或反击。

\subsection{第二条件：我的牌有合适的功能}

即使范围有优势，也不是所有牌都适合进攻。你还需要判断自己的具体手牌是否有合适功能。

适合进攻的牌通常包括：

\begin{itemize}
\item 强价值牌：希望更差牌跟注；
\item 坚果牌：希望构建大底池；
\item 强听牌：可以半诈唬；
\item 带阻断牌的空气牌：适合选择性诈唬；
\item 有后门权益的牌：被跟注后仍有改善空间；
\item 没有摊牌价值的牌：不下注几乎无法赢，只能选择诈唬或放弃。
\end{itemize}

不适合随意进攻的牌包括：

\begin{itemize}
\item 中等摊牌价值：经常更适合控池；
\item 弱摊牌价值：不值得投入太多筹码；
\item 低质量听牌：被跟注后权益不足；
\item 没有阻断、没有权益、没有故事的空气牌。
\end{itemize}

一手牌是否适合进攻，不只看它强不强，还要看它的功能。强牌负责价值，强听牌负责半诈唬，低摊牌价值空气牌可以选择性诈唬，中等牌力则经常负责防守和摊牌。

\subsection{第三条件：对手有足够多可以弃掉的牌}

进攻能否盈利，最终还要看对手是否会弃牌。

如果对手范围中有大量中等牌力、弱对子、错过听牌和空气牌，并且他有弃牌能力，那么你的进攻会产生弃牌权益。

如果对手范围非常强，或者对手类型是跟注站，那么进攻价值下降。此时你应该减少诈唬，增加价值下注。

判断对手是否适合被施压，可以问：

\begin{itemize}
\item 他翻牌是否跟得太宽？
\item 他转牌和河牌是否会放弃中等牌力？
\item 他是否害怕大底池？
\item 他是否会思考牌面变化？
\item 他是否有过明显 hero call 倾向？
\end{itemize}

对手如果没有弃牌按钮，你就不要用诈唬去寻找他的弃牌按钮。

\subsection{有效激进检查表}

在实战中，你可以使用以下检查表判断是否适合主动进攻：

\begin{enumerate}
\item 我是否拥有范围优势或坚果优势？
\item 我的具体手牌是价值牌、强听牌，还是好诈唬候选？
\item 这个牌面是否支持我的行动线？
\item 对手范围中是否有足够多中等牌力或弱牌？
\item 对手是否有弃牌能力？
\item 如果被跟注，我下一街是否有计划？
\item 如果被加注，我是否知道如何应对？
\end{enumerate}

如果大部分答案为“是”，进攻通常更有依据。
如果大部分答案为“否”，所谓进攻很可能只是鲁莽。

\section{牌例拆解}

\subsection{牌例一：高牌干燥面上的有效进攻}

\subsubsection*{牌局背景}

有效筹码 100BB。你在按钮位拿到 \hand{QJo}，前面玩家弃牌，你 open raise 到 2.5BB。大盲跟注。

翻牌为 \hand{A♠ 7♣ 2♦ rainbow}。大盲过牌，行动轮到你。

\subsubsection*{新手常见想法}

新手可能会想：

\begin{itemize}
\item “我没中牌，应该过牌。”
\item “如果他有 A，我下注就是送钱。”
\item “没有牌不要乱 bluff。”
\end{itemize}

这些想法只看到了自己的具体牌，没有看到范围优势和弃牌权益。

\subsubsection*{理性分析}

第一，范围优势。按钮 open range 中有大量 A、高对和强牌，而大盲防守范围较宽，包含许多没有击中的牌。\hand{A♠ 7♣ 2♦ rainbow} 是典型有利于翻前加注者的高牌干燥面。

第二，手牌功能。你持有 \hand{Q♠ J♣}，没有摊牌价值，属于空气牌。但在这个牌面上，它可以作为小尺度诈唬候选。它不适合摊牌赢，下注至少有机会让更好牌弃牌。

第三，对手弃牌范围。大盲可能持有 K 高、Q 高、小口袋对子、未击中的连张等。这些牌面对小注会有一部分弃牌。

第四，下注尺度。由于牌面干燥，小尺度下注即可发挥作用，不需要大额冒险。

\subsubsection*{推荐行动}

小尺度 C-bet，例如 1/3 底池。

\subsubsection*{本手牌教训}

没中牌不代表不能进攻。只要范围有优势、牌面适合、对手有可弃牌范围，空气牌也可以成为合理诈唬候选。

\subsection{牌例二：低连接面上不该盲目激进}

\subsubsection*{牌局背景}

有效筹码 100BB。你在前位拿到 \hand{AQo}，open raise 到 3BB。大盲跟注。

翻牌为 \hand{9♠ 8♠ 7♦ two-tone}。大盲过牌，行动轮到你。

\subsubsection*{新手常见想法}

一些新手会想：

\begin{itemize}
\item “我是翻前加注者，要持续下注。”
\item “如果我不打，就显得很弱。”
\item “他可能也没中，我可以代表大牌。”
\end{itemize}

这正是把主动权误用为自动进攻。

\subsubsection*{理性分析}

第一，范围优势。这个牌面非常容易击中大盲防守范围。大盲可以有 \hand{T♠ 8♣}、\hand{J♠ T♣}、\hand{9♠ 8♣}、\hand{8♠ 7♣}、\hand{7♠ 6♣}、小对子、两对、暗三条和各种听牌。你的前位 open range 虽然有高对，但在坚果优势上并不明显。

第二，手牌功能。你持有 \hand{AQo}，没有同花听牌，没有顺子听牌，只有两张高牌。它不是强听牌，也不是好的半诈唬候选。

第三，对手弃牌范围。大盲在这个牌面上有大量可继续牌。即使他没有成强牌，也可能有一对、顺子听牌、同花听牌或组合听牌。弃牌权益不足。

\subsubsection*{推荐行动}

倾向过牌。

这不是因为你永远不能下注，而是因为当前三个有效激进条件都不充分：范围优势不足，手牌功能较差，对手可弃牌范围有限。

\subsubsection*{本手牌教训}

主动权不是自动 C-bet 的理由。在低连接湿润面上，没有权益的高牌空气不适合盲目进攻。

\subsection{牌例三：强听牌的半诈唬加注}

\subsubsection*{牌局背景}

有效筹码 100BB。你在按钮位拿到 \hand{A♠ Q♠}，前面玩家弃牌，你 open raise 到 2.5BB。大盲跟注。

翻牌为 \hand{J♠ T♠ 4♦}。大盲主动下注 1/2 底池，行动轮到你。

\subsubsection*{新手常见想法}

新手可能会想：

\begin{itemize}
\item “我还没有成牌，先跟注看转牌。”
\item “加注太激进了，万一他有强牌怎么办？”
\item “等中了同花再打。”
\end{itemize}

跟注当然可以考虑，但这手牌也具备很强的半诈唬加注价值。

\subsubsection*{理性分析}

第一，手牌功能。你拥有坚果同花听牌、两张高牌和顺子相关可能，是非常强的听牌。

第二，弃牌权益。大盲主动下注范围中可能包含 Jx、Tx、弱听牌和部分空气牌。加注可以让其中一部分牌弃牌，尤其是没有能力承受大底池的中等牌力。

第三，被跟注后的权益。即使对手跟注，你仍然有许多转牌改善机会。任何黑桃、K，甚至 A 或 Q 都可能改善你的牌力或带来继续施压机会。

第四，范围表达。你在按钮位 open，可以拥有 \hand{J♠ J♣}、\hand{T♠ T♣}、\hand{J♠ T♣}、强同花听牌和强组合听牌。加注线并非不可信。

\subsubsection*{推荐行动}

可以加注半诈唬，也可以根据对手类型选择跟注。

面对能弃掉中等牌力的对手，加注更有吸引力。面对极度不弃牌但会支付完成听牌的对手，跟注实现权益也合理。

\subsubsection*{本手牌教训}

强听牌不只是用来被动追牌。它们可以通过半诈唬加注同时获得弃牌权益和牌面权益。

\subsection{牌例四：对跟注站不该强行 bluff}

\subsubsection*{牌局背景}

有效筹码 100BB。你在 cutoff 拿到 \hand{Q♠ T♠}，open raise 到 2.5BB。大盲是一名明显跟注站，喜欢用任何对子跟到河牌。

大盲跟注。

翻牌为 \hand{K♠ 7♣ 2♦ rainbow}。大盲过牌，你下注 1/3 底池，大盲跟注。

转牌为 \hand{4♠}。大盲过牌，行动轮到你。你没有成牌，也没有明显听牌。

\subsubsection*{新手常见想法}

一些新手会想：

\begin{itemize}
\item “我翻牌下注了，转牌要继续代表 K。”
\item “他可能只是 7 或小对子，我打大一点能打走。”
\item “如果我放弃，这手牌就输了。”
\end{itemize}

这些想法忽略了最关键的信息：对手类型。

\subsubsection*{理性分析}

第一，范围和牌面。你在 \hand{K♠ 7♣ 2♦} 上翻牌下注有一定合理性，因为高牌干燥面有利于进攻方。但转牌 \hand{4} 是空白牌，并没有显著增加你的故事力量。

第二，手牌功能。你没有成牌，也没有明显听牌或阻断效果。你的牌是低质量空气牌。

第三，对手弃牌能力。大盲是跟注站，翻牌跟注后很可能继续用 Kx、7x、小对子甚至 A 高跟注。你的弃牌权益较低。

第四，继续下注风险。转牌继续下注可能只是扩大亏损。如果他跟注，你河牌大概率仍然没有计划。

\subsubsection*{推荐行动}

过牌放弃。

\subsubsection*{本手牌教训}

对不会弃牌的人，少 bluff。激进必须建立在弃牌权益上。没有弃牌权益的进攻，只是把筹码送给对手。

\subsection{牌例五：价值加注的意义}

\subsubsection*{牌局背景}

有效筹码 100BB。你在大盲拿到 \hand{7♠ 7♣}。按钮位 open raise 到 2.5BB，你跟注。

翻牌为 \hand{K♠ 7♠ 2♦ two-tone}。你过牌，按钮下注 1/3 底池，行动轮到你。

\subsubsection*{新手常见想法}

一些新手会想：

\begin{itemize}
\item “我中了暗三条，跟注让他继续 bluff。”
\item “如果加注，他可能会弃牌。”
\item “慢打可以隐藏牌力。”
\end{itemize}

慢打可以考虑，但不应该成为默认动作。

\subsubsection*{理性分析}

第一，牌力功能。你持有暗三条，是非常强的价值牌，接近坚果。

第二，牌面结构。牌面有两张同花，存在同花听牌。按钮位可能有 Kx、同花听牌、口袋对子和部分空气牌。你有很多目标牌可以继续支付。

第三，加注目的。这是价值加注，同时带有保护意义。你希望 Kx、强听牌和部分不愿弃牌的手牌支付更多筹码。你也希望开始构建大底池，而不是只让对手用小注控制底池。

第四，底池规划。如果只跟注，转牌对手可能过牌，导致你很难在河牌赢到足够多筹码。翻牌加注可以更快建立底池。

\subsubsection*{推荐行动}

加注。

尺度可以根据对手和牌面选择，例如加到对手下注的 3 至 4 倍左右，目标是让 Kx 和听牌继续支付。

\subsubsection*{本手牌教训}

加注不只是 bluff，也可以是价值和构建底池。强牌不能总是慢打，尤其在有听牌和可支付目标时，价值加注非常重要。

\section{本章总结}

德州扑克不是被动等待好牌的游戏。如果你只靠摊牌赢底池，就永远少了一条腿。真正成熟的牌手，既能在摊牌时赢，也能在合适条件下通过进攻让对手弃牌。

本章最重要的概念是主动权和弃牌权益。主动权让你能够率先施压，弃牌权益让你即使当前不是最好牌，也有机会通过下注或加注直接赢下底池。

但激进不是鲁莽。有效激进必须满足三个条件：

\begin{center}
\textbf{我的范围有优势。\\
我的牌有合适的功能。\\
对手有足够多可以弃掉的牌。}
\end{center}

范围优势决定你的故事是否可信。手牌功能决定你的进攻是否有质量。对手可弃牌范围决定你的下注是否真的有弃牌权益。

适合激进的牌通常包括强价值牌、坚果牌、强听牌、带阻断效果或后门权益的诈唬候选。不适合盲目激进的牌包括中等摊牌价值、低质量听牌、没有权益且没有故事的空气牌。

适合进攻的牌面通常是有利于你范围的牌面，例如高牌干燥面、强化你范围的转牌，以及对手范围中有大量中等牌力的局面。不适合盲目进攻的牌面，则包括低连接湿润面、明显完成对手听牌的牌面和多人底池中的复杂牌面。

对手类型决定激进策略的盈利性。面对会弃牌的人，可以合理施压；面对跟注站，应该减少诈唬，增加价值下注；面对激进玩家，要避免用弱牌无计划对抗；面对紧弱玩家，可以增加偷池和施压频率。

加注是更强烈的进攻工具。它可以用于价值、保护、半诈唬和施压，但不能用于试探或情绪表达。加注之前，必须知道自己希望对手做什么，也必须知道面对再加注如何处理。

下一章我们将补充基础计算工具：outs、胜率与底池赔率。只有先理解“这个希望值不值这个价格”，才能在后续章节中更理性地决定何时跟注、何时弃牌。再往后，我们将进入弃牌与跟注，把数学计算与对手范围分析结合起来。

\section{课后训练}

\subsection*{训练一：有效激进三条件检查}

记录最近 10 手你主动下注或加注的牌，逐手回答：

\begin{enumerate}
\item 我的范围是否有优势？
\item 我的牌是否有合适功能？
\item 对手是否有足够多可以弃掉的牌？
\item 如果被跟注，我下一街是否有计划？
\item 如果被加注，我是否知道如何应对？
\end{enumerate}

如果三个核心条件中有两个以上不满足，说明这次进攻可能质量不高。

\subsection*{训练二：找出该进攻却被动的牌}

从最近牌局中找出 5 手你本可以主动进攻、但实际选择过牌或被动跟注的牌。

重点分析：

\begin{itemize}
\item 我是否拥有范围优势？
\item 我的牌是否是强听牌或好诈唬候选？
\item 对手是否容易弃牌？
\item 我当时为什么不敢下注或加注？
\item 如果重新来一次，我会如何设计进攻线？
\end{itemize}

这个训练的目标，是减少过度被动。

\subsection*{训练三：找出不该进攻却硬冲的牌}

从最近牌局中找出 5 手你进攻失败的牌，分析它们是否属于鲁莽进攻。

逐手回答：

\begin{enumerate}
\item 牌面是否真的有利于我的范围？
\item 我的牌是否有牌面权益、阻断效果或后门权益？
\item 对手是否具备弃牌能力？
\item 我是否因为情绪、追损或不想示弱而下注？
\item 我在哪一街应该停止进攻？
\end{enumerate}

这个训练的目标，是区分有效激进和鲁莽进攻。

\subsection*{训练四：半诈唬候选训练}

记录 10 手你持有听牌的牌，判断它们是否适合半诈唬。

重点回答：

\begin{itemize}
\item 这是强听牌还是弱听牌？
\item 被跟注后我有多少改善牌？
\item 我是否有弃牌权益？
\item 我是否有位置？
\item 对手是否会弃牌？
\item 加注、下注和跟注中哪一个更合适？
\end{itemize}

不要把所有听牌都当成半诈唬。强听牌才是高质量进攻候选。

\subsection*{训练五：对手施压对象分类}

在你的常见牌局中，选出 5 名对手，按照以下类型分类：

\begin{itemize}
\item 跟注站；
\item 紧弱玩家；
\item 激进玩家；
\item 松凶玩家；
\item 普通娱乐玩家。
\end{itemize}

对每位玩家回答：

\begin{enumerate}
\item 他是否容易弃牌？
\item 他在哪条街最容易放弃？
\item 他面对大注是否会过度弃牌？
\item 他是否喜欢抓诈唬？
\item 我应该对他增加价值下注，还是增加诈唬？
\end{enumerate}

这个训练的目标，是让激进策略和对手类型匹配。

\subsection*{训练六：加注目的复盘}

记录最近 10 次你加注的牌，逐手回答：

\begin{itemize}
\item 我是价值加注、保护加注、半诈唬加注，还是施压加注？
\item 哪些更差牌会继续？
\item 哪些更好牌会弃牌？
\item 如果对手再加注，我的计划是什么？
\item 这次加注是否让底池大小匹配我的牌力？
\end{itemize}

如果你无法说明加注目的，说明这次加注可能只是情绪化或试探性行动。

\chapter{基础计算：Outs、胜率与赔率}

\section{本章导入}

前面几章中，我们多次提到“听牌是否值得继续”“面对下注是否有足够胜率”“下注尺度会影响对手继续范围”“河牌跟注必须满足数学条件”。这些内容背后，都离不开一个最基础的问题：

\begin{center}
\textbf{继续这手牌，值不值当前这个价格？}
\end{center}

很多新手在牌桌上判断是否跟注时，依靠的是感觉。他们会想：“我还有同花机会”“我还有顺子机会”“他可能在 bluff”“都打到这里了，再看一张”。这些想法都有一个共同问题：它们只说明你还有机会，却没有说明这个机会是否值得你付出当前成本。

德州扑克不是只看有没有机会，而是看机会和价格是否匹配。
同花听牌当然有机会完成，但如果对手下注太大，你付出的价格可能不划算。
顶对当然有机会领先，但如果对手河牌大注范围几乎全是强牌，你的胜率可能不够。
一手牌当然有可能赢，但如果你长期赢不到足够比例，跟注就是亏损。

因此，在进入下一章“弃牌与跟注”之前，我们需要先补上一套基础计算工具。本章要讲三个最重要的入门概念：

\begin{itemize}
\item Outs：哪些未发出的牌能帮助你改善；
\item 胜率：你的牌最终赢下底池的大致概率；
\item 赔率：你为了继续游戏需要付出的价格是否划算。
\end{itemize}

本章不是数学课，也不要求你在牌桌上进行复杂精确计算。你只需要掌握几条简单但非常实用的规则：如何数 outs，如何用“2/4 法则”快速估算胜率，如何计算底池赔率，如何比较“实际胜率”和“所需胜率”。

本章的核心输出是“继续决策四步法”：

\begin{center}
\textbf{先数 outs。\\
再估胜率。\\
计算跟注所需胜率。\\
比较实际胜率与所需胜率。}
\end{center}

学完本章之后，你会明白一个最重要的原则：

\begin{center}
\textbf{能不能跟，不看“有没有希望”，而看“这个希望值不值这个价格”。}
\end{center}

\section{核心概念}

\subsection{什么是 Outs}

Outs 可以理解为：牌堆中还没有发出的牌里，哪些牌能让你的手牌改善为大概率领先的牌。中文里可以把它理解为“补牌”或“有效补牌”。

例如你手里是：

\begin{center}
\texttt{As 7s}
\end{center}

翻牌是：

\begin{center}
\texttt{Ks 9s 2d}
\end{center}

你现在有同花听牌。只要转牌或河牌再来一张黑桃，你就能组成同花。

一副牌中每种花色有 13 张牌。你已经看到 4 张黑桃：

\begin{center}
\texttt{As、7s、Ks、9s}
\end{center}

所以还剩：

\[
13 - 4 = 9
\]

张黑桃没有出现。这 9 张黑桃，就是你的 outs。

也就是说：

\begin{center}
\textbf{同花听牌通常有 9 个 outs。}
\end{center}

再看一个顺子听牌的例子。你手里是：

\begin{center}
\texttt{9c 8c}
\end{center}

翻牌是：

\begin{center}
\texttt{7d 6s Kh}
\end{center}

现在只要转牌或河牌出现 \texttt{5} 或 \texttt{T}，你就能组成顺子。一副牌中有 4 张 5 和 4 张 T，因此你一共有：

\[
4 + 4 = 8
\]

个 outs。

这就是两头顺听牌，通常有 8 个 outs。

Outs 的意义在于，它把“我还有机会”变成了更具体的数字。你不再只是模糊地说“我有听牌”，而是能够说清楚“我有 9 个 outs”或“我有 8 个 outs”。

\subsection{常见 Outs 类型}

新手可以先记住几类最常见的 outs。

\begin{center}
\begin{tabular}{lll}
\toprule
\textbf{听牌类型} & \textbf{典型情况} & \textbf{常见 outs} \\
\midrule
同花听牌 & 四张同花，差一张成同花 & 9 \\
两头顺听牌 & 两端都能成顺 & 8 \\
卡顺听牌 & 只有中间一张能成顺 & 4 \\
两张高牌成一对 & 两张高牌都可能配对 & 约 6 \\
一对变三条 & 口袋对子或一对想中三条 & 2 \\
同花听牌加两头顺 & 组合听牌 & 约 15 \\
\bottomrule
\end{tabular}
\end{center}

例如：

\begin{itemize}
\item 同花听牌：通常 9 outs；
\item 两头顺听牌：通常 8 outs；
\item 卡顺听牌：通常 4 outs；
\item 两张高牌想击中一对：通常最多 6 outs；
\item 一对想变三条：通常 2 outs；
\item 同花听牌加两头顺：可能有 15 outs 左右。
\end{itemize}

这里要注意，“常见 outs”只是入门估算。实战中，你还要判断这些 outs 是否真的有效。有些牌虽然能让你成牌，但可能也会让对手成更强牌。这种 outs 不能全部算满。

\subsection{干净 Outs 与不干净 Outs}

不是所有 outs 都是干净的。

干净 outs 指的是：如果这张牌出现，你大概率能成为最好牌。
不干净 outs 指的是：这张牌虽然改善了你，但也可能让你输给更强牌。

例如你手里是：

\begin{center}
\texttt{As 7s}
\end{center}

翻牌是：

\begin{center}
\texttt{Ks 9s 2d}
\end{center}

你追的是 A 高同花。如果后面来黑桃，你大概率能组成很强的同花。因为你手中有 \texttt{As}，你的同花是坚果同花。这类 outs 比较干净。

但如果你手里是：

\begin{center}
\texttt{8s 7s}
\end{center}

翻牌是：

\begin{center}
\texttt{Ks Qs 2d}
\end{center}

你同样有同花听牌，但即使后面来黑桃，你也只是 8 高同花。如果对手持有 \texttt{AsJs}、\texttt{AsTs}、\texttt{JsTs} 等更高同花，你完成同花后仍然可能输掉大底池。

这就是不干净 outs。

再看顺子听牌。你如果追的是低端顺子，而牌面可能同时让对手组成更高顺子，那么部分 outs 也要打折。比如你成了小顺子，但对手可能成更大顺子，这就是反向风险。

因此，新手不能机械地认为“我有 9 个 outs，所以一定可以跟”。你还要问：

\begin{itemize}
\item 我完成后是否接近坚果？
\item 我的 outs 是否可能让对手也完成更强牌？
\item 我的同花是否可能被更高同花压制？
\item 我的顺子是否可能输给更大顺子？
\item 我的高牌配对后是否可能输给更好踢脚？
\end{itemize}

如果 outs 不干净，就要适当打折。你可以不必精确计算成 7.5 个 outs 或 6.2 个 outs，但必须知道：低质量听牌的真实价值低于表面数字。

\subsection{什么是胜率}

胜率指的是：从当前局面到摊牌，你最终赢下底池的概率。

在听牌场景中，我们通常先用 outs 来估算胜率。例如你有 9 个 outs 的同花听牌，就可以估算你在转牌或河牌完成同花的大致概率。

新手不需要一开始掌握复杂概率公式，只需要先记住一个非常实用的近似方法：

\begin{center}
\textbf{翻牌到河牌：outs $\times 4 \approx$ 成牌概率。\\
转牌到河牌：outs $\times 2 \approx$ 成牌概率。}
\end{center}

这就是常说的“2/4 法则”。

例如你在翻牌有同花听牌，通常是 9 个 outs。翻牌之后还有转牌和河牌两张牌可以发，因此：

\[
9 \times 4 = 36\%
\]

也就是说，翻牌同花听牌到河牌完成同花的概率大约是 36\%。实际精确值约为 35\%，这个估算已经足够实战使用。

如果到了转牌，你仍然只有同花听牌，只剩河牌一张牌：

\[
9 \times 2 = 18\%
\]

实际精确值约为 19.6\%。同样，这个估算已经足够帮助你做初步决策。

\subsection{2/4 法则}

2/4 法则是新手最实用的胜率估算方法。

\begin{center}
\begin{tabular}{lll}
\toprule
\textbf{当前街道} & \textbf{还剩公共牌} & \textbf{估算方法} \\
\midrule
翻牌圈 & 转牌和河牌两张 & outs $\times 4$ \\
转牌圈 & 河牌一张 & outs $\times 2$ \\
\bottomrule
\end{tabular}
\end{center}

例如：

\begin{itemize}
\item 同花听牌 9 outs：翻牌约 36\%，转牌约 18\%；
\item 两头顺听牌 8 outs：翻牌约 32\%，转牌约 16\%；
\item 卡顺听牌 4 outs：翻牌约 16\%，转牌约 8\%；
\item 一对变三条 2 outs：翻牌约 8\%，转牌约 4\%；
\item 强组合听牌 15 outs：翻牌约 60\%，转牌约 30\%。
\end{itemize}

这条规则有两个重要作用。

第一，它让你快速知道听牌大概有多少胜率。你不再只是说“我有同花机会”，而是能说“我大约有 36\% 的完成概率”。

第二，它让你能和底池赔率进行比较。你可以判断当前跟注价格是否划算。

需要注意的是，2/4 法则估算的是“改善成目标牌”的概率，而不一定等于“最终赢牌概率”。因为你完成听牌后仍可能输给更强牌，你没完成听牌时也可能偶尔靠高牌或对子赢。因此它是实用估算，而不是绝对答案。

\subsection{什么是底池赔率}

底池赔率回答的问题是：

\begin{center}
\textbf{面对当前下注，我跟注需要赢多少比例，才不亏？}
\end{center}

计算公式是：

\[
\text{所需胜率} =
\frac{\text{跟注金额}}{\text{跟注后总底池}}
\]

这里要特别注意，是“跟注后总底池”，不是当前底池。

例如，当前底池为 100，对手下注 50。你需要跟注 50。若你跟注，底池中会有：

\[
100 + 50 + 50 = 200
\]

其中包括原底池 100，对手下注 50，以及你的跟注 50。

因此，你需要的胜率是：

\[
\frac{50}{200} = 25\%
\]

也就是说，只要你长期能赢超过 25\%，这个跟注才可能盈利。如果你的实际胜率低于 25\%，跟注就是亏损。

再看一个例子。当前底池为 100，对手下注 100。你要跟注 100。跟注后总底池为：

\[
100 + 100 + 100 = 300
\]

所需胜率为：

\[
\frac{100}{300} = 33.3\%
\]

也就是说，面对一个底池大小下注，你大约需要 33.3\% 的胜率。

底池赔率的核心不是让你背公式，而是让你建立价格意识：

\begin{center}
\textbf{小注需要的胜率低，大注需要的胜率高。}
\end{center}

\subsection{常见下注尺度对应的所需胜率}

新手可以先记住几个常见下注尺度对应的所需胜率。

\begin{center}
\begin{tabular}{llll}
\toprule
\textbf{对手下注} & \textbf{原底池} & \textbf{跟注后总底池} & \textbf{所需胜率} \\
\midrule
1/3 底池 & 100 & 166.7 & 约 20\% \\
1/2 底池 & 100 & 200 & 25\% \\
2/3 底池 & 100 & 233.3 & 约 28.6\% \\
1 倍底池 & 100 & 300 & 33.3\% \\
1.5 倍底池 & 100 & 400 & 37.5\% \\
\bottomrule
\end{tabular}
\end{center}

这张表非常重要。它会直接影响你面对下注时的防守范围。

面对 1/3 底池小注，你只需要约 20\% 胜率，因此可以用较宽范围继续。
面对 1/2 底池下注，你需要 25\% 胜率。
面对底池大小下注，你需要约 33.3\% 胜率。
面对超池下注，你需要的胜率更高，必须明显收紧跟注范围。

很多新手的问题是，不管对手下注多大，他们都用差不多的牌去跟。这是严重错误。下注尺度越大，你继续所需的胜率越高，对手范围通常也越强，你必须更加谨慎。

\subsection{需要胜率与实际胜率}

跟注决策中，有两个概念必须区分：需要胜率和实际胜率。

需要胜率，是由底池赔率决定的。
实际胜率，是你的手牌面对对手范围时，实际能赢的概率。

比如当前底池 100，对手下注 50。你跟注需要 25\% 胜率。这是需要胜率。

但你到底有没有 25\% 胜率，要看你的牌面对手下注范围的表现。如果你有同花听牌，可能有大约 18\% 到 36\% 的改善概率，具体看是在翻牌还是转牌。如果你在河牌拿着中等对子面对对手大注，你可能只能击败诈唬，实际胜率取决于对手诈唬比例。

很多新手只看需要胜率，却不判断实际胜率。例如他们会说：“这个下注不大，我赔率不错，可以跟。”但如果对手几乎全是强价值牌，你实际胜率远远不够，再好的价格也不能乱跟。

也有相反情况。对手下注较大，你所需胜率较高，但如果你有强组合听牌，实际胜率也可能足够，甚至可以主动加注半诈唬。

所以，正确的比较是：

\begin{center}
\textbf{实际胜率 $\geq$ 所需胜率：可以考虑继续。\\
实际胜率 $<$ 所需胜率：通常应该弃牌。}
\end{center}

这里说“可以考虑继续”，不是说一定要跟。因为实战还要考虑位置、对手类型、隐含赔率、反向隐含赔率和后续行动。

\subsection{隐含赔率}

隐含赔率指的是：你现在跟注虽然从直接底池赔率看不一定完全划算，但如果后续成牌，可能还能从对手身上赢到更多筹码。

最典型例子是小对子追暗三条。

你翻前用 \texttt{55} 跟注对手加注。你现在的目标不是靠 5 高对子摊牌赢，而是希望翻牌击中暗三条。如果你中了暗三条，而对手持有 \texttt{AA}、\texttt{KK} 或强顶对，他可能会支付很多筹码。你当前跟注的价值，部分来自未来可能赢到的大底池，这就是隐含赔率。

同花听牌和顺子听牌也可能有隐含赔率。你现在跟注追听牌，如果后面完成强牌，对手愿意继续支付，你的跟注价值就会上升。

但隐含赔率成立需要条件：

\begin{itemize}
\item 有效筹码足够深；
\item 你完成后牌力足够强；
\item 对手会在你成牌后继续支付；
\item 你的成牌不明显，不容易让对手立刻停止下注；
\item 你有位置，方便控制后续行动。
\end{itemize}

如果筹码很浅，即使你成牌也赢不到更多，隐含赔率就低。
如果你追的是明显同花，而第三张同花一出现对手立刻不再支付，隐含赔率也低。
如果你追的是低同花，完成后可能输给更高同花，隐含赔率看似高，实则包含反向风险。

因此，隐含赔率不是“我中了可以赢很多”的幻想，而是要判断对手是否真的会在你成牌后付钱。

\subsection{反向隐含赔率}

反向隐含赔率是新手非常容易忽略的概念。它指的是：你成牌后看似变强了，但实际上可能输给更强牌，结果反而输掉更大的底池。

典型例子包括：

\begin{itemize}
\item 弱 A 击中 A，但输给更强 A；
\item 低同花完成同花，但输给更高同花；
\item 小顺子完成顺子，但输给更大顺子；
\item 非坚果听牌完成后，被对手更强听牌压制；
\item 两张高牌配对后，输给更好踢脚。
\end{itemize}

例如你持有：

\begin{center}
\texttt{Ac 6d}
\end{center}

翻牌为：

\begin{center}
\texttt{Ah Js 8c}
\end{center}

你击中了顶对，但踢脚很弱。如果对手范围中有很多 \texttt{AK}、\texttt{AQ}、\texttt{AJ}、\texttt{AT}，你虽然“中了 A”，却很容易输给更强 A。此时你的牌有很强的反向隐含赔率。

再如你持有：

\begin{center}
\texttt{8s 7s}
\end{center}

牌面为：

\begin{center}
\texttt{Ks Qs 2d}
\end{center}

你追同花，但即使完成，也可能输给更高同花。你看似有 9 个 outs，但这些 outs 并不完全干净。

反向隐含赔率提醒我们：

\begin{center}
\textbf{不是所有“成牌”都值得追，不是所有“中了”都舒服。}
\end{center}

很多边缘牌和弱听牌的问题，不是它们完全没有机会，而是它们在成牌后仍然可能输大底池。

\subsection{把 Outs、胜率和赔率连起来}

真正实战中，outs、胜率和赔率不是孤立使用的。它们应该连成一个完整决策流程。

\begin{enumerate}
\item 先数 outs：哪些牌能帮助我改善？
\item 再估胜率：用 2/4 法则估算大致完成概率；
\item 计算所需胜率：根据底池赔率判断我需要赢多少；
\item 比较两者：实际胜率是否足够？
\item 调整判断：outs 是否干净？是否有隐含赔率？是否有反向隐含赔率？
\end{enumerate}

例如，你翻牌有同花听牌，9 outs。用 2/4 法则估算，翻牌到河牌约有 36\% 完成概率。如果对手下注给你的价格只需要 25\% 胜率，那么从直接赔率看，可以考虑继续。

但如果这是低同花听牌，且对手范围中可能有更高同花，你要给 outs 打折。
如果对手很强，成牌后不一定支付，你的隐含赔率下降。
如果你没有位置，转牌还可能面对更大下注，你也要谨慎。

所以，入门计算不是让你机械跟注，而是让你有一个理性起点。你先用数字判断价格是否大致合理，再结合牌面、范围和对手类型做最终决策。

\section{新手常见误区}

\subsection{误区一：只要有 Outs 就继续}

很多新手看到自己有听牌，就觉得必须跟注。他们会说：“我还有 9 张牌能中同花。”“我还有 8 张牌能成顺。”但他们没有问：对手下注给我的价格是否划算？

有 outs 只是说明你有机会，并不说明这个机会值得购买。面对小注，同花听牌可能轻松继续；面对转牌大注，同样的同花听牌可能就不够价格。

正确思路是：先数 outs，再看价格。没有价格意识的追牌，是长期亏损的重要来源。

\subsection{误区二：把所有 Outs 都算满}

新手常常机械地数 outs，却不判断 outs 是否干净。

低同花听牌不应该和坚果同花听牌等价。
弱 A 击中 A 不应该和强 A 等价。
小顺子听牌不应该和坚果顺子听牌等价。

如果你的 outs 可能让你成第二好牌，就要打折。很多时候，问题不是你中不了，而是你中了以后还会输更多。

\subsection{误区三：翻牌胜率和转牌胜率混淆}

翻牌听牌和转牌听牌差别很大。

翻牌同花听牌还有两张牌，约 36\%。
转牌同花听牌只剩一张牌，约 18\%。

很多新手翻牌跟注合理，但转牌面对大注仍然机械继续，就是因为他们没有意识到胜率已经明显下降。到了转牌，听牌继续条件更加苛刻。

\subsection{误区四：只看赔率，不看对手范围}

有些玩家学会底池赔率后，又走向另一个误区：只要价格看起来不错，就跟注。

但底池赔率只告诉你“需要多少胜率”，不告诉你“实际有没有这么多胜率”。实际胜率必须结合对手范围。

例如河牌面对小注，你可能只需要较低胜率。但如果对手是极度被动玩家，河牌下注几乎全是价值牌，你的中等牌力仍然可能不能跟。

价格重要，范围同样重要。

\subsection{误区五：幻想隐含赔率}

很多新手会说：“我中了就能赢很多。”但这常常只是幻想。

如果对手看到同花完成就不再支付，你的同花听牌没有你想象中那么高隐含赔率。
如果有效筹码很浅，你即使成牌也赢不到更多。
如果你追的是低同花，完成后可能输给更高同花，那么所谓隐含赔率可能变成反向隐含赔率。

隐含赔率必须建立在真实支付可能性上，而不是想象中的大底池。

\subsection{误区六：忽略位置和后续行动}

听牌跟注不仅要看当前赔率，还要看后续行动。

你在有位置时跟注，后续可以根据对手行动决定是否继续。
你在无位置时跟注，转牌可能再次面对下注，无法轻松实现权益。

尤其在翻牌面对下注时，如果你跟注后转牌还可能被大额施压，那么你不能只看当前一街价格。你要考虑整手牌路径。

\section{实战思考框架}

本章的核心框架是“继续决策四步法”：

\begin{center}
\textbf{先数 outs。\\
再估胜率。\\
计算跟注所需胜率。\\
比较实际胜率与所需胜率。}
\end{center}

\subsection{第一步：先数 Outs}

面对下注时，如果你还没有成强牌，而是依靠后续改善继续，第一步就是数 outs。

你要问：

\begin{itemize}
\item 哪些牌能让我成同花？
\item 哪些牌能让我成顺子？
\item 哪些牌能让我从一对变三条？
\item 哪些高牌配对后可能让我领先？
\item 这些 outs 是否真的干净？
\end{itemize}

不要只说“我有听牌”。要尽量具体：我是 9 outs 同花听牌，还是 4 outs 卡顺？我是坚果同花听牌，还是低同花听牌？这些差别会直接影响后续决策。

\subsection{第二步：再估胜率}

数出 outs 后，用 2/4 法则估算胜率。

\begin{itemize}
\item 翻牌到河牌：outs $\times 4$；
\item 转牌到河牌：outs $\times 2$。
\end{itemize}

例如：

\begin{itemize}
\item 翻牌 9 outs：约 36\%；
\item 转牌 9 outs：约 18\%；
\item 翻牌 8 outs：约 32\%；
\item 转牌 8 outs：约 16\%；
\item 翻牌 4 outs：约 16\%；
\item 转牌 4 outs：约 8\%。
\end{itemize}

这里的胜率是估算值，不是绝对精确数字。它的作用是帮助你快速判断“我大概有多少机会”。

\subsection{第三步：计算跟注所需胜率}

然后，根据底池赔率计算你需要多少胜率。

公式是：

\[
\text{所需胜率} =
\frac{\text{跟注金额}}{\text{跟注后总底池}}
\]

如果不想每次精确计算，可以先记住常见尺度：

\begin{center}
\begin{tabular}{ll}
\toprule
\textbf{下注尺度} & \textbf{所需胜率} \\
\midrule
1/3 底池 & 约 20\% \\
1/2 底池 & 25\% \\
2/3 底池 & 约 28.6\% \\
1 倍底池 & 33.3\% \\
1.5 倍底池 & 37.5\% \\
\bottomrule
\end{tabular}
\end{center}

下注越大，你需要的胜率越高。面对大注，用弱听牌继续通常会变得很困难。

\subsection{第四步：比较实际胜率与所需胜率}

最后，把估算胜率和所需胜率进行比较。

如果你的实际胜率高于所需胜率，可以考虑跟注。
如果你的实际胜率低于所需胜率，通常应该弃牌。

例如你翻牌有同花听牌，约 36\% 完成概率。对手下注给你的价格只需要 25\% 胜率，从直接赔率看可以跟注。

但如果你转牌仍然只有同花听牌，约 18\% 完成概率，对手下注底池大小，你需要约 33.3\% 胜率。此时如果没有很强隐含赔率或额外胜率，跟注通常不划算。

\subsection{第五步：用质量因素修正判断}

前四步给出的是基础数学判断。实战中还要修正。

你需要继续问：

\begin{itemize}
\item 我的 outs 是否干净？
\item 我完成后是否接近坚果？
\item 对手是否会在我成牌后继续支付？
\item 我是否有位置？
\item 有效筹码是否足够深？
\item 如果我没中，后续是否还能 bluff？
\item 如果我中了，是否可能输给更强牌？
\end{itemize}

这些因素不会推翻基础计算，但会影响最终决策。强听牌、有位置、深筹码、对手愿意支付，会让继续更有吸引力。弱听牌、无位置、浅筹码、反向隐含赔率高，则会让继续变差。

\section{牌例拆解}

\subsection{牌例一：翻牌同花听牌面对半池下注}

\subsubsection*{牌局背景}

有效筹码 100BB。你在按钮位拿到：

\begin{center}
\texttt{As 7s}
\end{center}

你 open raise 到 2.5BB，大盲跟注。

翻牌为：

\begin{center}
\texttt{Ks 9s 2d}
\end{center}

大盲下注 1/2 底池，行动轮到你。

\subsubsection*{新手常见想法}

新手可能会想：

\begin{itemize}
\item “我有同花听牌，应该跟。”
\item “如果来黑桃，我就是同花。”
\item “A 高同花应该很强。”
\end{itemize}

这些想法方向不错，但还需要用计算确认价格。

\subsubsection*{理性分析}

第一，数 outs。你有 A 高同花听牌。一共有 13 张黑桃，你已经看到 4 张黑桃，因此还有 9 张黑桃可以帮助你成同花。

第二，估胜率。翻牌到河牌还有两张牌，用 2/4 法则：

\[
9 \times 4 = 36\%
\]

你完成同花的概率大约是 36\%。

第三，计算所需胜率。对手下注 1/2 底池。面对半池下注，你大约需要 25\% 胜率。

第四，比较。你的估算完成概率约 36\%，高于所需胜率 25\%。而且你是 A 高同花听牌，outs 较干净。

\subsubsection*{推荐行动}

可以跟注，也可以根据对手类型和弃牌权益考虑加注半诈唬。

如果对手容易弃牌，加注有吸引力；如果对手不爱弃牌，跟注实现权益也合理。

\subsubsection*{本手牌教训}

翻牌坚果同花听牌面对半池下注，通常有足够直接赔率继续。关键不是“有听牌就跟”，而是胜率确实高于所需胜率。

\subsection{牌例二：转牌同花听牌面对底池下注}

\subsubsection*{牌局背景}

有效筹码 100BB。你在按钮位拿到：

\begin{center}
\texttt{As 7s}
\end{center}

翻牌为：

\begin{center}
\texttt{Ks 9s 2d}
\end{center}

翻牌你面对小注跟注。

转牌为：

\begin{center}
\texttt{4h}
\end{center}

大盲下注底池大小，行动轮到你。

\subsubsection*{新手常见想法}

很多新手会想：

\begin{itemize}
\item “我还是同花听牌，继续跟。”
\item “河牌来黑桃就赢。”
\item “翻牌都跟了，转牌不能弃。”
\end{itemize}

这正是常见的转牌追牌错误。

\subsubsection*{理性分析}

第一，数 outs。你仍然有 9 个同花 outs。

第二，估胜率。现在只剩河牌一张牌，用 2/4 法则：

\[
9 \times 2 = 18\%
\]

你完成同花的概率大约是 18\%。

第三，计算所需胜率。面对底池大小下注，你大约需要 33.3\% 胜率。

第四，比较。18\% 明显低于 33.3\%。从直接赔率看，跟注不划算。

第五，考虑隐含赔率。你是 A 高同花听牌，完成后很强。但如果河牌出现第三张黑桃，对手是否还会继续支付很多？如果不会，你的隐含赔率不足以弥补当前价格。

\subsubsection*{推荐行动}

多数情况下弃牌。

除非对手极有可能在河牌继续支付大量筹码，或者你有额外胜率和明确后续计划，否则面对转牌底池下注不应机械追听牌。

\subsubsection*{本手牌教训}

同样是同花听牌，翻牌和转牌价值完全不同。翻牌有两张牌可看，转牌只剩一张。转牌面对大注时，听牌继续条件更苛刻。

\subsection{牌例三：两头顺听牌面对小注}

\subsubsection*{牌局背景}

有效筹码 100BB。你在大盲拿到：

\begin{center}
\texttt{9c 8c}
\end{center}

按钮位 open raise 到 2.5BB，你跟注。

翻牌为：

\begin{center}
\texttt{7d 6s Kh}
\end{center}

你过牌，按钮下注 1/3 底池。

\subsubsection*{新手常见想法}

新手可能会想：

\begin{itemize}
\item “我还没成牌，但有顺子听牌。”
\item “下注不大，可以看一张。”
\item “来 5 或 T 就好了。”
\end{itemize}

这手牌可以用计算确认。

\subsubsection*{理性分析}

第一，数 outs。你需要 \texttt{5} 或 \texttt{T} 成顺。4 张 5 加 4 张 T，共 8 outs。

第二，估胜率。翻牌到河牌：

\[
8 \times 4 = 32\%
\]

第三，计算所需胜率。面对 1/3 底池下注，你大约需要 20\% 胜率。

第四，比较。32\% 高于 20\%。从直接赔率看，继续合理。

第五，考虑其他因素。你在大盲无位置，这是不利因素；但下注较小，听牌较清晰，继续可以接受。如果对手容易弃牌，也可以考虑加注半诈唬；如果对手不爱弃牌，跟注实现权益即可。

\subsubsection*{推荐行动}

可以跟注。根据对手类型，也可以考虑加注半诈唬。

\subsubsection*{本手牌教训}

两头顺听牌面对小注通常有足够价格继续。但无位置会降低实现权益的能力，因此不能只看 outs，还要考虑后续行动。

\subsection{牌例四：卡顺听牌面对大注}

\subsubsection*{牌局背景}

有效筹码 100BB。你在按钮位拿到：

\begin{center}
\texttt{Qc Jc}
\end{center}

你 open raise，大盲跟注。

翻牌为：

\begin{center}
\texttt{Ad Ks 4h}
\end{center}

大盲过牌，你下注 1/3 底池，大盲 check-raise 到接近底池大小加注。

\subsubsection*{新手常见想法}

新手可能会想：

\begin{itemize}
\item “我来 T 就是顺子。”
\item “我有卡顺，不能马上弃。”
\item “如果中了，可能赢大底池。”
\end{itemize}

但卡顺听牌只有 4 个 outs，面对大注必须非常谨慎。

\subsubsection*{理性分析}

第一，数 outs。你只有 \texttt{T} 能成顺，一共 4 张 T，因此是 4 outs。

第二，估胜率。翻牌到河牌：

\[
4 \times 4 = 16\%
\]

第三，计算所需胜率。面对接近底池大小的加注，你通常需要接近或超过 33\% 的胜率。

第四，比较。16\% 明显不够。

第五，考虑隐含赔率。虽然成顺后很强，但对手 check-raise 表示强范围，后续你未必中牌后能轻松拿到足够价值。此外，如果转牌没中，你很可能还要面对更大下注。

\subsubsection*{推荐行动}

弃牌。

\subsubsection*{本手牌教训}

卡顺听牌只有 4 outs，不能因为“中了很强”就面对大注继续。低 outs 听牌需要非常好的价格才值得追。

\subsection{牌例五：小对子追暗三条}

\subsubsection*{牌局背景}

有效筹码 150BB。你在按钮位拿到：

\begin{center}
\texttt{5c 5d}
\end{center}

前位玩家 open raise 到 3BB。你在按钮位行动。

\subsubsection*{新手常见想法}

新手可能会想：

\begin{itemize}
\item “我有对子，可以跟。”
\item “如果翻牌来 5，就很隐蔽。”
\item “不中就弃，也不会亏太多。”
\end{itemize}

这类手牌的价值主要来自隐含赔率。

\subsubsection*{理性分析}

第一，数 outs。如果你想从口袋对子变暗三条，牌堆中还剩 2 张 5，因此只有 2 outs。

第二，翻牌击中暗三条的概率并不高。粗略理解即可：大约八九次中一次。也就是说，大多数时候你会错过翻牌。

第三，为什么仍然可以跟？因为有效筹码深，你有位置，而且如果你中了暗三条，对手前位强范围可能有 \texttt{AA}、\texttt{KK}、\texttt{AK}、强顶对等愿意支付。这提供了隐含赔率。

第四，什么时候不能跟？如果有效筹码很浅，你即使中了也赢不到更多；如果后面还有激进玩家可能 squeeze；如果对手很紧且不中牌不会支付，那么跟注价值下降。

\subsubsection*{推荐行动}

在有效筹码较深、你有位置、后面再加注风险不高时，可以跟注。若筹码较浅或后面玩家激进，应更倾向弃牌。

\subsubsection*{本手牌教训}

小对子追暗三条不是靠直接赔率，而是靠隐含赔率。你不是因为经常中，而是因为一旦中牌可能赢大底池。

\subsection{牌例六：低同花听牌的反向隐含赔率}

\subsubsection*{牌局背景}

有效筹码 150BB。你在大盲拿到：

\begin{center}
\texttt{8s 6s}
\end{center}

按钮位 open raise，你跟注。

翻牌为：

\begin{center}
\texttt{Ks Qs 2d}
\end{center}

你过牌，按钮下注 2/3 底池。

\subsubsection*{新手常见想法}

新手可能会想：

\begin{itemize}
\item “我有同花听牌，9 outs。”
\item “如果来黑桃，我就是同花。”
\item “筹码深，中了可以赢很多。”
\end{itemize}

但这手牌有明显反向隐含赔率。

\subsubsection*{理性分析}

第一，表面 outs。你看似有 9 个黑桃 outs。

第二，outs 质量。你追的是 8 高同花。如果按钮有 \texttt{AsXs}、\texttt{JsTs}、\texttt{Ts9s} 等更高同花听牌，你完成同花后可能输大底池。

第三，位置问题。你在大盲无位置。即使转牌来黑桃，你也很难控制底池，并且对手如果强势下注，你会很尴尬。

第四，下注尺度。面对 2/3 底池下注，需要约 28.6\% 胜率。翻牌到河牌表面有约 36\%，但考虑 outs 不干净、无位置和反向隐含赔率，真实价值下降。

\subsubsection*{推荐行动}

可以谨慎跟注，但不能把它当作强听牌。面对更大下注、激进对手或后续难实现权益时，弃牌也完全合理。

\subsubsection*{本手牌教训}

低同花听牌的表面 outs 很好看，但并不等于真实价值很高。成牌后可能输给更高同花，是反向隐含赔率的典型来源。

\section{本章总结}

本章补充了德州扑克新手必须掌握的基础计算工具：outs、胜率和赔率。

Outs 是还能帮助你改善的补牌。常见同花听牌有 9 outs，两头顺听牌有 8 outs，卡顺听牌有 4 outs，一对变三条有 2 outs。数 outs 是判断听牌价值的第一步。

胜率是你最终赢下底池的大致概率。新手最实用的方法是 2/4 法则：

\begin{center}
\textbf{翻牌到河牌：outs $\times 4$。\\
转牌到河牌：outs $\times 2$。}
\end{center}

底池赔率告诉你，面对当前下注，你跟注需要赢多少比例才不亏。公式是：

\[
\text{所需胜率} =
\frac{\text{跟注金额}}{\text{跟注后总底池}}
\]

小注需要的胜率低，大注需要的胜率高。面对 1/3 底池下注约需 20\% 胜率，面对半池下注约需 25\%，面对底池下注约需 33.3\%。

真正的跟注判断，需要比较实际胜率和所需胜率：

\begin{center}
\textbf{实际胜率高于所需胜率，可以考虑继续。\\
实际胜率低于所需胜率，通常应该弃牌。}
\end{center}

但计算只是起点。实战中，你还必须考虑 outs 是否干净、是否有隐含赔率、是否存在反向隐含赔率、是否有位置、对手是否会支付、后续是否容易实现权益。

本章最重要的一句话是：

\begin{center}
\textbf{能不能跟，不看“有没有希望”，而看“这个希望值不值这个价格”。}
\end{center}

下一章“弃牌与跟注”将建立在本章基础之上，进一步讨论面对下注时如何估算对手价值牌和诈唬牌，哪些牌适合 bluff-catch，哪些看似强的牌应该弃，以及如何克服“不想被诈唬”的心理。

\section{课后训练}

\subsection*{训练一：Outs 速算}

找出 20 个翻牌或转牌听牌局面，逐一写下：

\begin{enumerate}
\item 我是什么类型的听牌？
\item 我有多少 outs？
\item 这些 outs 是否干净？
\item 如果不干净，原因是什么？
\end{enumerate}

训练目标是让你从“我有听牌”升级为“我有多少有效补牌”。

\subsection*{训练二：2/4 法则练习}

对以下常见 outs 进行胜率估算：

\begin{itemize}
\item 9 outs；
\item 8 outs；
\item 6 outs；
\item 4 outs；
\item 2 outs；
\item 15 outs。
\end{itemize}

分别计算：

\begin{enumerate}
\item 翻牌到河牌的大致概率；
\item 转牌到河牌的大致概率。
\end{enumerate}

训练目标是形成快速估算直觉。

\subsection*{训练三：底池赔率计算}

假设当前底池为 100，分别面对以下下注，计算你需要多少胜率：

\begin{itemize}
\item 对手下注 25；
\item 对手下注 33；
\item 对手下注 50；
\item 对手下注 75；
\item 对手下注 100；
\item 对手下注 150。
\end{itemize}

使用公式：

\[
\text{所需胜率} =
\frac{\text{跟注金额}}{\text{跟注后总底池}}
\]

训练目标不是追求小数点精确，而是熟悉下注越大、所需胜率越高。

\subsection*{训练四：听牌是否值得跟}

记录最近 10 手你用听牌跟注的手牌，逐手回答：

\begin{enumerate}
\item 我有多少 outs？
\item 我在翻牌还是转牌？
\item 我的估算胜率是多少？
\item 对手下注给我的所需胜率是多少？
\item 实际胜率是否高于所需胜率？
\item 我是否考虑了隐含赔率和反向隐含赔率？
\end{enumerate}

训练目标是减少无脑追听牌。

\subsection*{训练五：隐含赔率判断}

找出 5 手你追小对子、同花听牌或顺子听牌的手牌，分析：

\begin{itemize}
\item 有效筹码是否足够深？
\item 如果我成牌，对手是否会支付？
\item 我的成牌是否足够隐蔽？
\item 我是否有位置？
\item 如果这些条件不满足，我是否高估了隐含赔率？
\end{itemize}

训练目标是避免用“中了赢很多”作为模糊理由。

\subsection*{训练六：反向隐含赔率复盘}

找出 5 手你成牌后仍然输大底池的手牌，分析是否存在反向隐含赔率。

重点回答：

\begin{enumerate}
\item 我追的是不是非坚果牌？
\item 我完成后是否可能输给更强同花、更大顺子或更好踢脚？
\item 我是否把不干净 outs 算满了？
\item 我是否在成牌后过度自信？
\item 如果重来一次，我会不会在前面街更谨慎？
\end{enumerate}

训练目标是理解：不是所有成牌都值得追，也不是所有成牌后都适合打大底池。

\chapter{弃牌与跟注}

\section{本章导入}

这是新手最重要、也最反直觉的一章。上一章我们学习了 outs、胜率与底池赔率的基础计算；本章在此基础上，进一步讨论面对下注时如何估算对手范围、选择 bluff-catcher，以及如何用弃牌纪律保护长期 EV。

很多娱乐玩家以为，自己最大的漏洞是不够激进、不敢诈唬、不懂高级打法。事实上，在大多数低级别牌局和娱乐局中，玩家最大的亏损来源往往不是诈唬太少，而是跟注太多。

他们翻前跟注太多，翻牌跟注太多，转牌跟注太多，河牌面对大注仍然跟注太多。只要手里有一对、有听牌、有高牌、有一点希望，就舍不得弃牌。他们常常不是因为理性判断而跟注，而是因为不甘心、好奇、害怕被诈唬、已经投入太多筹码，或者单纯想“看一眼”。

但德州扑克中，跟注不是情绪行为，而是数学行为。每一次跟注，本质上都是在用一定成本购买摊牌或继续实现权益的机会。你必须判断：我需要赢多少比例？我的实际胜率是否足够？对手下注范围中有多少价值牌？有多少诈唬牌？我的牌能击败哪些部分？

如果这些问题回答不清，却只是因为“他可能在诈唬”而跟注，那么长期结果通常会非常糟糕。

弃牌则恰恰相反。很多新手把弃牌看成软弱，认为弃牌就是被欺负、被打跑、错过机会。但成熟牌手知道，弃牌是一种纪律。它不是承认失败，而是承认当前价格、范围和牌力不支持继续。能在错误价格下放弃，能在大底池中放下第二好牌，能克服“不想被诈唬”的心理，是长期盈利的重要能力。

本章的核心输出是“跟注四问”：

\begin{center}
\textbf{我需要多少胜率？\\
对手有哪些价值牌？\\
对手有哪些诈唬牌？\\
我的牌能否击败足够多的下注范围？}
\end{center}

只要你能在面对下注时稳定问这四个问题，就会开始从“好奇型跟注”转向“理性防守”。

\section{核心概念}

\subsection{为什么过度跟注是最大亏损来源}

过度跟注之所以危险，是因为它看起来不像错误。

一次糟糕的诈唬失败，损失很明显；一次错误的大额下注，也容易被记住。但跟注不同。跟注经常披着“合理怀疑”的外衣：我有一对，我有听牌，我赔率好像还行，他也可能在诈唬，我已经投了很多。于是玩家一手又一手地跟，一条街又一条街地跟，最终长期损失大量筹码。

过度跟注会带来几个问题。

第一，你会用太弱的范围进入后续街。翻牌应该弃掉的牌，被你带到了转牌；转牌应该弃掉的牌，被你带到了河牌；河牌应该弃掉的牌，又被你付钱看摊牌。每多跟一条街，错误成本都会增加。

第二，你会让对手的价值下注变得非常赚钱。如果你总是用中等牌力和弱牌支付，对手根本不需要复杂诈唬，只需要不断用强牌下注，就能从你身上拿价值。

第三，你会强化自己的坏习惯。偶尔一次抓到诈唬，会让你记住“我跟对了”；但更多时候你支付给对手价值牌的筹码，却容易被归因于“运气不好”“他刚好有”。这种记忆偏差会让过度跟注持续存在。

第四，你会失去资金和情绪稳定性。河牌错误跟注往往发生在大底池中，一次错误可能抵消很多小底池盈利。连续几次错误跟注后，牌手容易进入追损和上头状态。

娱乐玩家最大的漏洞，常常不是他们完全不会进攻，而是他们无法正确防守。面对下注时，他们跟得太宽、弃得太少、想得太情绪化。要想进步，必须先学会控制跟注频率。

\subsection{跟注不是情绪行为，而是数学行为}

跟注的本质，是用一定成本换取赢下当前底池的机会。它应该是一个数学和范围判断问题，而不是情绪问题。

情绪化跟注常见于以下心理：

\begin{itemize}
\item “我不想被诈唬。”
\item “我已经投了这么多，现在不能弃。”
\item “他不可能每次都有。”
\item “我想看他到底是什么。”
\item “这手牌弃了太弱。”
\item “我今天已经被打走太多次了。”
\end{itemize}

这些想法都不是合格的跟注理由。真正的跟注理由应该是：

\begin{itemize}
\item 当前底池赔率要求我只需要赢一定比例；
\item 我的牌能击败对手下注范围中足够多的组合；
\item 对手有足够多诈唬；
\item 我的牌在防守范围中处于较高位置；
\item 我持有合适阻断牌，减少对手强价值组合；
\item 对手下注尺度和行动线不支持过强价值范围。
\end{itemize}

跟注是一项需要证明的行动。你不能因为“可能赢”就跟注，因为任何牌都有可能赢。你需要判断的是：这个可能性是否足以覆盖跟注成本。

如果你无法说明自己为什么能赢足够多，弃牌通常比跟注更好。

\subsection{底池赔率的基本概念}

底池赔率，是判断跟注是否合理的基础工具。它回答的问题是：

\begin{center}
\textbf{面对当前下注，我需要赢多少比例，跟注才不亏？}
\end{center}

假设底池中已有 100，面对对手下注 50。此时你需要跟注 50，跟注后总底池为 200。你投入 50 去争夺 200 的总底池，因此你需要至少赢：

\[
\frac{50}{200} = 25\%
\]

也就是说，只要你长期能赢超过 25\%，这个跟注在数学上才可能盈利。

再看几个常见例子：

\begin{center}
\begin{tabular}{llll}
\toprule
\textbf{对手下注} & \textbf{原底池} & \textbf{跟注后总底池} & \textbf{所需胜率} \\
\midrule
1/3 底池 & 100 & 166.7 & 约 20\% \\
1/2 底池 & 100 & 200 & 25\% \\
2/3 底池 & 100 & 233.3 & 约 28.6\% \\
1 倍底池 & 100 & 300 & 33.3\% \\
超池 1.5 倍 & 100 & 400 & 37.5\% \\
\bottomrule
\end{tabular}
\end{center}

这个表说明，下注越大，你跟注需要的胜率越高。面对小注，你可以用更宽范围继续；面对大注，你必须显著收紧跟注范围。

但是，底池赔率只是第一步。知道自己需要 25\% 胜率，并不代表你真的有 25\% 胜率。你还需要估算对手下注范围中价值牌和诈唬牌的比例。

\subsection{需要胜率与实际胜率}

跟注决策中，必须区分两个概念：需要胜率和实际胜率。

需要胜率由底池赔率决定。面对 1/2 底池下注，你大约需要 25\% 胜率；面对底池大小下注，你需要约 33.3\% 胜率。

实际胜率则取决于你的牌面对手下注范围的表现。也就是说，你的牌能击败对手多少价值牌和诈唬牌。

很多新手的错误在于，只看到底池很大或下注不算太大，就觉得“赔率不错，可以跟”。但他们没有认真判断自己的实际胜率是否达到要求。

例如河牌底池 100，对手下注 50，你需要约 25\% 胜率。看起来门槛不高。但如果对手下注范围中 90\% 是价值牌，10\% 是诈唬，而你的牌只能击败诈唬，那么你的实际胜率只有 10\%，跟注就是亏损。

反过来，如果对手有大量错过听牌，价值牌组合并不多，你只需要 25\% 胜率，而你的牌能击败 35\% 的下注范围，那么跟注就是合理的。

所以，跟注不能只看价格，也不能只看牌力。必须把价格和范围结合起来：

\begin{center}
\textbf{需要胜率由下注尺度决定；实际胜率由对手范围决定。}
\end{center}

只有实际胜率大于需要胜率，跟注才有意义。

\subsection{面对下注时如何估算对手范围}

面对下注时，新手最常问的是：“他有没有？”但这个问题太模糊。正确问题应该是：

\begin{center}
\textbf{他用哪些价值牌这样下注？他用哪些诈唬牌这样下注？}
\end{center}

估算对手范围，可以按以下步骤进行。

第一，回顾翻前行动。对手是前位加注者、大盲防守者、冷跟注者，还是 3-bet 方？不同翻前行动决定了他的初始范围。

第二，回顾每条街行动。他翻牌下注还是过牌？转牌是否继续？河牌突然下注还是连续开火？每一次行动都会过滤掉一部分牌。

第三，结合牌面变化。哪些听牌完成了？哪些听牌错过了？哪些转牌或河牌强化了对手范围？哪些牌对他的故事不自然？

第四，列出价值牌。包括强顶对、两对、暗三条、顺子、同花、葫芦等。要尽量具体，而不是笼统地说“他可能有强牌”。

第五，列出诈唬牌。包括错过的同花听牌、错过的顺子听牌、没有摊牌价值的浮动牌、带阻断牌的空气牌等。关键是判断这些牌是否真的会用当前尺度下注。

第六，结合对手类型。被动玩家的下注通常偏价值；激进玩家的诈唬更多；跟注站很少在河牌主动 bluff；紧弱玩家大注通常很强。

对手范围不是凭空想象，而是由翻前范围、牌面结构、行动线和对手类型共同决定。

\subsection{哪些牌适合 bluff-catch}

Bluff-catch，即抓诈唬，指的是一手牌通常打不过对手的价值下注，但可以击败对手诈唬。抓诈唬牌的价值不在于它能打赢强牌，而在于对手诈唬足够多时，它可以通过跟注盈利。

适合 bluff-catch 的牌通常具备以下特点：

\begin{itemize}
\item 有一定摊牌价值，可以击败对手空气牌；
\item 不足以主动价值下注；
\item 面对对手价值牌大多落后；
\item 在自己的防守范围中不算太弱；
\item 持有合适阻断牌，减少对手强价值组合；
\item 不阻断对手可能的诈唬组合。
\end{itemize}

例如河牌牌面完成同花，对手大注。你持有一张关键花色 A，但没有成同花。这个 A 可能阻断对手坚果同花组合，使你的抓诈唬价值上升。
又如对手可能错过同花听牌，你手中没有该花色牌，反而不阻断他的错过听牌诈唬组合，这有时也会让跟注更合理。

但新手不要过度迷信 blocker。阻断牌只是辅助因素。最关键的仍然是：对手是否真的有足够诈唬。

很多娱乐玩家河牌大注极少 bluff。面对这类对手，再漂亮的 bluff-catcher 也应该谨慎跟注。

\subsection{哪些牌看似强但应该弃}

新手最难接受的是：有些看似强的牌，在特定行动线和牌面下应该弃。

例如顶对。顶对在翻牌可能很强，但如果对手连续三条街大注，牌面又完成了顺子或同花，你的顶对可能只剩抓诈唬价值。它看起来是顶对，实际上已经无法击败对手大多数价值牌。

例如超对。\hand{A♠ A♣} 在翻前极强，但在 \hand{9♠ 8♣ 7♦ 6♥ T♠} 这样的公共牌上，面对对手持续强力进攻，可能只是非常脆弱的一对。

例如两对。两对在干燥牌面上很强，但在四连顺或三同花牌面上，面对大注可能不再适合支付。

例如小同花。你完成了同花，但如果牌面允许更高同花，对手又表现出极强行动，小同花可能会被更大同花支配。

是否弃牌，不取决于牌名看起来强不强，而取决于它能击败对手当前下注范围的多少部分。

新手需要接受一个事实：

\begin{center}
\textbf{一手牌的绝对牌力很强，不代表它在当前局面中仍然强。}
\end{center}

很多大亏损，正是来自不愿意放下“看起来很强但实际上只是在抓诈唬”的牌。

\subsection{河牌跟注的苛刻条件}

河牌跟注是最苛刻的跟注，因为河牌没有未来权益。你不能靠后续补牌改善，不能再实现听牌，不能寄希望于下一张牌。你跟注后，结果立即揭晓。

因此，河牌跟注必须满足更严格的条件。

第一，你必须知道自己需要多少胜率。下注越大，需要胜率越高。

第二，你必须列出对手的价值牌。如果对手有大量合理价值牌，而你几乎都打不过，那么跟注门槛很高。

第三，你必须列出对手的诈唬牌。如果你找不到足够自然的错过听牌或空气牌，那就不要幻想对手有很多 bluff。

第四，你必须判断对手是否真的会诈唬。很多被动娱乐玩家河牌大注极少诈唬。面对他们，应该过度弃牌，而不是过度跟注。

第五，你的牌必须在防守范围中足够靠前。如果你用所有 bluff-catcher 都跟注，就一定跟太多。你应该选择更好的 bluff-catcher，而不是随便用任何一对支付。

河牌跟注前，最危险的句子是：“我想看一眼。”

看一眼不是免费的。河牌看一眼，往往要付出一大注。你必须确保这次“看一眼”有足够数学依据。

\subsection{如何克服“不想被诈唬”的心理}

“不想被诈唬”是新手过度跟注的核心心理之一。

很多玩家宁愿付钱看错，也不愿弃牌后怀疑自己被 bluff。他们觉得被诈唬是一种羞辱，觉得弃牌就是软弱，觉得对手可能会因此继续欺负自己。于是他们用跟注来维护自尊，而不是维护 EV。

成熟牌手必须理解：被诈唬不是失败，错误跟注才是失败。

如果对手用诈唬赢了你一手小底池，但你长期在不合适的价格下弃牌，你仍然可能是盈利的。反过来，如果你为了不被诈唬而在河牌频繁跟注，你会给对手的价值牌支付大量筹码。

克服这种心理，可以从三个原则开始。

第一，允许自己被诈唬。没有人能抓住所有 bluff。只要你的整体防守频率合理，被个别诈唬打走并不可怕。

第二，关注长期结果，而不是单手自尊。你不是为了证明自己聪明而跟注，而是为了做长期盈利决策。

第三，用范围和价格替代情绪。每次想“我不信”时，改问：“我需要多少胜率？他的诈唬够不够？”

弃牌不是被欺负。弃牌是拒绝在不利价格下支付。

\section{新手常见误区}

\subsection{误区一：已经投了钱，所以不能弃}

这是典型的沉没成本错误。很多新手会想：“我翻牌跟了，转牌也跟了，现在河牌弃掉太可惜。”但已经投入底池的筹码不再属于你。当前决策只应该考虑：现在跟注是否盈利。

如果河牌下注让你需要 33\% 胜率，而你判断自己只有 10\% 胜率，那么即使你前面已经投入很多，也应该弃牌。继续跟注不会把前面的筹码救回来，只会增加新的亏损。

\subsection{误区二：只要对手可能诈唬，就跟注}

对手当然可能诈唬。问题不是“有没有可能”，而是“够不够多”。

如果你面对底池下注，需要约 33\% 胜率。对手下注范围中如果只有 10\% 是诈唬，你跟注就是亏损。即使这一次他刚好在诈唬，也不能证明你的跟注长期正确。

跟注需要的是足够多的诈唬，而不是理论上的可能性。

\subsection{误区三：不看下注尺度}

有些新手面对不同尺度的下注，跟注范围几乎一样。面对 1/3 底池也跟，面对底池下注也跟，面对超池下注还想跟。这是非常严重的错误。

下注越大，你需要胜率越高，对手范围通常也越极化。你必须随着尺度增大而收紧跟注范围。

小注可以用较多边缘牌防守；大注必须用更强的牌和更好的 bluff-catcher 防守。

\subsection{误区四：只看自己的牌名}

“我有顶对。”“我有两对。”“我有同花。”“我有 AA。”这些描述都不够。

跟注时要看的是：我的牌能击败对手下注范围中的哪些牌？如果只能击败诈唬，那它就是 bluff-catcher；如果对手诈唬不足，看起来再强也应该弃。

牌名不能替代范围判断。

\subsection{误区五：跟注为了看信息}

有些玩家河牌跟注的理由是：“我想看看他到底是什么牌。”但信息不是免费的。尤其在河牌，大注看信息的成本非常高。

如果你总是为了信息而跟注，对手只需要用强牌下注，就能从你身上稳定赚钱。你可以通过观察行动线、摊牌记录、对手习惯来获取信息，而不是用错误跟注购买昂贵信息。

\subsection{误区六：不会选择抓诈唬牌}

新手经常随便用一对抓诈唬，却不知道哪些牌更适合 bluff-catch。好的抓诈唬牌通常应当阻断对手价值牌，且不阻断对手诈唬牌。

例如对手可能用错过同花听牌 bluff，如果你手里拿着该花色牌，可能反而减少了他错过同花听牌的组合。相反，如果你没有阻断这些错过听牌，但阻断了他部分强价值牌，跟注质量可能更高。

当然，选择 bluff-catcher 的前提仍然是：对手真的有足够诈唬。

\section{实战思考框架}

本章的核心框架是“跟注四问”：

\begin{center}
\textbf{我需要多少胜率？\\
对手有哪些价值牌？\\
对手有哪些诈唬牌？\\
我的牌能否击败足够多的下注范围？}
\end{center}

\subsection{第一问：我需要多少胜率？}

面对下注，第一步先看价格。你要知道自己跟注需要赢多少比例。

可以先记住几个常见数字：

\begin{center}
\begin{tabular}{ll}
\toprule
\textbf{对手下注尺度} & \textbf{你大约需要的胜率} \\
\midrule
1/3 底池 & 20\% \\
1/2 底池 & 25\% \\
2/3 底池 & 28.6\% \\
1 倍底池 & 33.3\% \\
1.5 倍底池 & 37.5\% \\
\bottomrule
\end{tabular}
\end{center}

这些数字不需要在实战中精确到小数点，但你必须有大致概念。面对越大的下注，你需要越强的理由继续。

\subsection{第二问：对手有哪些价值牌？}

接下来，列出对手可能的价值牌。

不要笼统地说“他可能有强牌”，而要具体：

\begin{itemize}
\item 他能有哪些两对？
\item 他能有哪些暗三条？
\item 他能有哪些顺子？
\item 他能有哪些同花？
\item 他会不会用顶对这样下注？
\item 他前面的行动是否支持这些强牌？
\end{itemize}

列出价值牌后，你要判断自己的牌能击败多少。如果对手价值牌大量存在，而你的牌几乎都打不过，那么跟注必须依赖对手有足够诈唬。

\subsection{第三问：对手有哪些诈唬牌？}

然后，列出对手可能的诈唬牌。

常见诈唬来源包括：

\begin{itemize}
\item 错过的同花听牌；
\item 错过的顺子听牌；
\item 翻牌浮动、转牌获得权益但河牌错过的牌；
\item 持有阻断牌但没有摊牌价值的牌；
\item 由于牌面变化而适合代表强牌的空气牌。
\end{itemize}

你需要判断这些牌是否真实存在于对手范围中。比如一个被动玩家翻牌跟注、转牌跟注、河牌突然大注，他真的会用错过听牌这样打吗？如果不会，你不能强行给他添加大量诈唬。

\subsection{第四问：我的牌能否击败足够多的下注范围？}

最后，把前三个问题结合起来。

如果你需要 25\% 胜率，而对手下注范围中至少有 30\% 是你能击败的牌，跟注可能合理。

如果你需要 33\% 胜率，但对手几乎全是价值牌，你只能击败极少数诈唬，那么应该弃牌。

如果你的牌不仅能击败诈唬，还能击败对手部分薄价值下注，那么跟注质量更高。
如果你的牌只能击败诈唬，那你必须确认对手诈唬足够多。

这就是理性跟注的完整逻辑。

\subsection{弃牌纪律检查表}

当你犹豫是否弃牌时，可以用以下问题提醒自己：

\begin{enumerate}
\item 我是否只是因为不甘心而想跟？
\item 我是否只是因为已经投入筹码而不想弃？
\item 我是否能列出对手足够多的诈唬？
\item 我的牌是否只能击败诈唬？
\item 对手类型是否支持他这样 bluff？
\item 如果我弃牌，即使这次被 bluff，长期是否仍然合理？
\end{enumerate}

如果你发现跟注理由主要来自情绪，而不是范围和价格，弃牌通常是更好的选择。

\section{牌例拆解}

\subsection{牌例一：河牌顶对面对被动玩家大注}

\subsubsection*{牌局背景}

有效筹码 100BB。你在按钮位拿到 \hand{KQo}，open raise 到 2.5BB。大盲是一名偏被动的娱乐玩家，平时很少主动诈唬。

大盲跟注。

翻牌为 \hand{Q♠ 8♣ 2♦ rainbow}。大盲过牌，你下注 1/3 底池，大盲跟注。

转牌为 \hand{5♠}。大盲过牌，你下注 2/3 底池，大盲跟注。

河牌为 \hand{J♠}。大盲突然下注接近底池。

\subsubsection*{新手常见想法}

很多新手会想：

\begin{itemize}
\item “我有顶对好踢脚，不能轻易弃。”
\item “他可能错过听牌在诈唬。”
\item “我已经打了两条街，现在弃太可惜。”
\item “万一他看我害怕河牌 J bluff 呢？”
\end{itemize}

这些想法都需要经过跟注四问检验。

\subsubsection*{理性分析}

第一，我需要多少胜率？面对接近底池下注，你大约需要 33\% 左右的胜率。

第二，对手有哪些价值牌？大盲可能有 \hand{Q♠ J♣} 河牌两对，\hand{8♠ 8♣}、\hand{5♠ 5♣}、\hand{2♠ 2♣} 暗三条，也可能有慢打强牌。被动玩家河牌突然主动大注，通常价值倾向很强。

第三，对手有哪些诈唬牌？翻牌较干燥，转牌也没有明显前门同花听牌。自然错过的听牌并不多。即使有一些 \hand{T♠ 9♣}、\hand{9♠ 7♣} 类牌，也要问：这个偏被动玩家真的会用它们河牌大注吗？答案通常是否定的。

第四，我的牌能否击败足够多下注范围？你的 \hand{K♠ Q♣} 能击败部分诈唬和少数过度下注的弱 Q，但面对对手价值范围大多落后。如果诈唬不足，实际胜率达不到跟注要求。

\subsubsection*{推荐行动}

弃牌。

\subsubsection*{本手牌教训}

顶对好踢脚不是自动跟注牌。面对被动玩家河牌突然大注，必须尊重其价值范围。弃掉顶对不是软弱，而是纪律。

\subsection{牌例二：转牌听牌面对大注}

\subsubsection*{牌局背景}

有效筹码 100BB。你在按钮位拿到 \hand{Q♠ J♠}，open raise 到 2.5BB。大盲跟注。

翻牌为 \hand{A♠ 7♠ 2♦}。大盲过牌，你下注 1/3 底池，大盲跟注。

转牌为 \hand{4♥}。大盲突然下注接近底池。

\subsubsection*{新手常见想法}

新手可能会想：

\begin{itemize}
\item “我有同花听牌，不能弃。”
\item “河牌来黑桃就能赢大底池。”
\item “我翻牌都打了，现在继续很正常。”
\end{itemize}

这手牌的关键，是转牌只剩一张牌，面对大注需要严格计算。

\subsubsection*{理性分析}

第一，我需要多少胜率？面对接近底池下注，你需要约 33\% 胜率。

第二，我的实际权益如何？你有同花听牌，但只剩一张河牌。普通同花听牌完成概率不足以单独支撑面对接近底池的大注。你可能还有 Q、J 高张价值，但面对大盲主动大注，这些高张未必有效。

第三，对手有哪些价值牌？大盲可能有强 A、两对、暗三条，甚至一些转牌改善牌。

第四，对手有哪些诈唬牌？他可能有部分弱听牌或空气牌，但需要判断其是否会主动大注。若对手不是激进玩家，诈唬比例未必足够。

第五，隐含赔率是否足够？即使河牌完成同花，你是否一定能再赢很多？如果大盲看到第三张黑桃后会停止支付，你的隐含赔率就不高。

\subsubsection*{推荐行动}

多数情况下弃牌，除非对手非常激进、下注范围中有大量诈唬，或你有明确隐含赔率。

\subsubsection*{本手牌教训}

有听牌不等于必须跟。转牌听牌面对大注时，只剩一张牌，继续条件比翻牌苛刻得多。

\subsection{牌例三：合格的河牌 bluff-catch}

\subsubsection*{牌局背景}

有效筹码 100BB。你在大盲拿到 \hand{A♥ Q♦}。按钮位 open raise 到 2.5BB，你跟注。

翻牌为 \hand{Q♥ 7♠ 2♥}。你过牌，按钮下注 1/3 底池，你跟注。

转牌为 \hand{4♣}。你过牌，按钮下注 2/3 底池，你跟注。

河牌为 \hand{8♦}。前门红心同花没有完成。你过牌，按钮下注 3/4 底池。

按钮是一名较激进玩家，过去多次在错过听牌时开火。

\subsubsection*{新手常见想法}

新手可能会想：

\begin{itemize}
\item “我有顶对好踢脚，应该跟。”
\item “他可能在 bluff。”
\item “这手牌不能弃。”
\end{itemize}

这次跟注可能合理，但仍然需要完整分析。

\subsubsection*{理性分析}

第一，我需要多少胜率？面对 3/4 底池下注，你需要约 30\% 左右胜率。

第二，对手有哪些价值牌？按钮可能有 \hand{A♠ Q♣}、\hand{K♠ Q♣}、\hand{Q♠ J♣}、\hand{7♠ 7♣}、\hand{2♠ 2♣}、\hand{4♠ 4♣}、少量两对。你的 \hand{A♠ Q♣} 能击败部分较弱 Q，输给暗三条和两对。

第三，对手有哪些诈唬牌？前门红心同花错过。按钮作为激进玩家，可能用错过红心听牌、\hand{J♠ T♣}、\hand{T♠ 9♣}、\hand{6♠ 5♣} 等牌继续 bluff。

第四，阻断牌如何？你持有 \hand{A♥}，阻断了对手部分坚果红心听牌诈唬，也阻断部分强 A 高红心组合。这一点有利有弊。它减少对手一些错过红心 bluff，但也减少对手部分 \hand{A♥ Q♥} 类强价值或强听牌组合。具体要结合对手范围判断。

第五，我的牌能否击败足够多下注范围？你有顶对好踢脚，能击败部分较弱 Q 和诈唬。面对激进玩家，这可能达到跟注要求。

\subsubsection*{推荐行动}

可以跟注。

\subsubsection*{本手牌教训}

抓诈唬不是因为“不想弃”，而是因为对手有足够诈唬、你的牌在防守范围中足够强，并且当前价格支持跟注。

\subsection{牌例四：看似强的同花也可能该弃}

\subsubsection*{牌局背景}

有效筹码 150BB。你在 cutoff 拿到 \hand{8♠ 7♠}，open raise 到 2.5BB。按钮跟注，大盲跟注。

翻牌为 \hand{K♠ T♠ 2♦}。大盲过牌，你下注 1/2 底池，按钮跟注，大盲弃牌。

转牌为 \hand{3♠}，你完成同花。你下注 2/3 底池，按钮跟注。

河牌为 \hand{Q♦}。你下注 1/2 底池，按钮大额加注接近全下。

按钮是一名偏紧被动玩家，很少在河牌大额诈唬。

\subsubsection*{新手常见想法}

很多新手会想：

\begin{itemize}
\item “我已经成同花了，怎么能弃？”
\item “他可能有两对或暗三条乱加注。”
\item “我等了这么久中花，必须看。”
\end{itemize}

但这是一个需要非常谨慎的局面。

\subsubsection*{理性分析}

第一，我需要多少胜率？面对接近全下的大额加注，你需要较高胜率。

第二，对手有哪些价值牌？按钮冷跟注范围中可以有许多更高同花，尤其是 \hand{A♠ X♠}、\hand{Q♠ J♠}、\hand{J♠ 9♠}、\hand{Q♠ 9♠} 等。转牌完成同花后，他跟注；河牌面对你下注大额加注，价值范围极强。

第三，对手有哪些诈唬牌？紧被动玩家河牌大额加注 bluff 非常少。公共牌已经有三张黑桃，他用两对或暗三条转诈唬的可能性较低。

第四，我的牌能击败哪些下注范围？你只有 8 高同花，能击败的价值同花很少，主要只能击败诈唬。但对手诈唬不足。

\subsubsection*{推荐行动}

弃牌。

\subsubsection*{本手牌教训}

完成同花不等于一定能打大底池。低同花面对紧被动玩家河牌大额加注，常常只是 bluff-catcher，甚至是很差的 bluff-catcher。

\subsection{牌例五：不要因为已经投入而河牌跟注}

\subsubsection*{牌局背景}

有效筹码 100BB。你在中位拿到 \hand{A♠ A♣}，open raise 到 3BB。大盲跟注。

翻牌为 \hand{J♠ T♠ 9♦ two-tone}。大盲过牌，你下注 2/3 底池，大盲跟注。

转牌为 \hand{8♠}。大盲过牌，你过牌。

河牌为 \hand{7♠}。大盲下注底池大小。

\subsubsection*{新手常见想法}

新手可能会想：

\begin{itemize}
\item “我有 AA，不能轻易弃。”
\item “我翻牌已经下注了。”
\item “他可能错过同花听牌。”
\item “底池这么大，弃了太可惜。”
\end{itemize}

这里最需要克服的，就是沉没成本和牌名执念。

\subsubsection*{理性分析}

第一，我需要多少胜率？面对底池下注，你需要约 33\% 胜率。

第二，对手有哪些价值牌？公共牌 \hand{J♠ T♣ 9♦ 8♥ 7♠} 极度连接。大盲可以有大量顺子组合，许多两对和暗三条也可能转化为强牌。你的 \hand{A♠ A♣} 在这个牌面上只是单对。

第三，对手有哪些诈唬牌？可能有错过同花听牌，但在如此连接的牌面上，对手有大量自然价值牌。并且大盲底池下注通常表示较强范围。

第四，我的牌能否击败足够多下注范围？你主要只能击败诈唬和少数过度下注的弱牌。若对手不是高频诈唬玩家，跟注不合理。

\subsubsection*{推荐行动}

弃牌。

\subsubsection*{本手牌教训}

\hand{A♠ A♣} 只是翻前最强起手牌，不是所有公共牌上的最强牌。已经投入的筹码不能成为河牌继续支付的理由。

\section{本章总结}

本章讨论的是新手最重要、也最反直觉的能力：弃牌与跟注。

娱乐玩家最大的漏洞，通常不是诈唬太少，而是跟注太多。他们害怕被诈唬，不愿放弃已经投入的筹码，想看对手到底有什么，于是不断用不够强的范围支付。长期来看，这会让对手的价值下注变得极其赚钱。

跟注不是情绪行为，而是数学行为。每一次跟注都必须回答四个问题：

\begin{center}
\textbf{我需要多少胜率？\\
对手有哪些价值牌？\\
对手有哪些诈唬牌？\\
我的牌能否击败足够多的下注范围？}
\end{center}

需要胜率由底池赔率决定，实际胜率由对手范围决定。只有当你的实际胜率大于需要胜率时，跟注才有意义。

河牌跟注尤其苛刻，因为河牌没有未来权益。你不能再靠补牌改善，不能再“看一张”，只能用当前牌力面对对手下注范围。如果你的牌只能击败诈唬，那么你必须确认对手真的有足够多诈唬。

弃牌不是软弱，而是纪律。被诈唬一手牌并不可怕，长期错误跟注才可怕。优秀牌手允许自己被合理诈唬，但不会为了维护自尊，在明显不利的价格下不断支付。

到这里，本书基础部分的核心逻辑已经完成：翻前建立范围，翻牌理解结构，转牌重新评估，河牌做最终判断；再进一步，你需要根据牌力功能决定任务，根据下注目的使用筹码，根据激进条件主动争取底池，并在面对进攻时用理性框架决定跟注或弃牌。

真正的进步，不是学会某一个固定答案，而是建立一套可重复、可复盘、可修正的决策系统。

\section{课后训练}

\subsection*{训练一：跟注四问记录}

下一次打牌后，记录 10 手你面对下注选择跟注的牌。每一手回答：

\begin{enumerate}
\item 我当时需要多少胜率？
\item 对手有哪些价值牌？
\item 对手有哪些诈唬牌？
\item 我的牌能否击败足够多的下注范围？
\item 我的跟注是理性判断，还是情绪驱动？
\end{enumerate}

训练目标是让每一次跟注都有明确理由。

\subsection*{训练二：河牌错误跟注复盘}

找出最近 5 手你河牌跟注后输牌的手牌，逐手分析：

\begin{itemize}
\item 我是否只是在抓诈唬？
\item 对手是否真的有足够诈唬？
\item 我是否忽略了下注尺度？
\item 我是否因为不想被诈唬而跟注？
\item 如果重来一次，我是否应该弃牌？
\end{itemize}

这个训练的目标，是减少河牌高成本错误。

\subsection*{训练三：成功弃牌记录}

记录 5 手你成功弃牌的手牌，哪怕你不知道对手最终是什么牌。

重点写下：

\begin{enumerate}
\item 我为什么选择弃牌？
\item 对手有哪些价值牌？
\item 对手诈唬是否不足？
\item 我的牌是否只是 bluff-catcher？
\item 即使这次被诈唬，长期弃牌是否仍然合理？
\end{enumerate}

这个训练的目标，是强化“弃牌也是正确决策”的正反馈。

\subsection*{训练四：底池赔率练习}

随机选择 20 个下注场景，计算你需要的大致胜率：

\begin{itemize}
\item 对手下注 1/3 底池；
\item 对手下注 1/2 底池；
\item 对手下注 2/3 底池；
\item 对手下注 1 倍底池；
\item 对手下注 1.5 倍底池。
\end{itemize}

不要求精确到小数点，但要形成尺度越大、跟注门槛越高的直觉。

\subsection*{训练五：bluff-catcher 筛选}

找出 10 手你在河牌持有一对或中等牌力面对下注的牌，判断它们是否是好的 bluff-catcher。

逐手回答：

\begin{itemize}
\item 我的牌能否击败对手部分价值下注？
\item 如果不能，它是否只能击败诈唬？
\item 我是否阻断对手强价值牌？
\item 我是否阻断对手诈唬牌？
\item 对手类型是否支持他有足够诈唬？
\end{itemize}

训练目标是学会选择高质量跟注牌，而不是随便用任何一对抓诈唬。

\subsection*{训练六：克服不想被诈唬}

每次你因为“不想被诈唬”而想跟注时，暂停并写下：

\begin{center}
\textbf{被诈唬一手牌不可怕，长期错误跟注才可怕。}
\end{center}

然后回答：

\begin{enumerate}
\item 我是否能列出足够多诈唬组合？
\item 这个对手是否真的会 bluff？
\item 我跟注是为了盈利，还是为了看牌？
\item 如果我弃牌，即使这次被 bluff，长期是否仍然合理？
\end{enumerate}

这个训练的目标，是把跟注从自尊反应，转化为理性决策。


\backmatter

\nocite{*}
\printbibliography[heading=bibintoc,title={参考文献}]

\end{document}
